\documentclass[12pt,a4paper,openany]{book}

% Подключение русского языка и кодировки
\usepackage[utf8]{inputenc}
\usepackage[T2A]{fontenc}
\usepackage[russian]{babel}
\usepackage{tikz}
\usepackage{tikz-3dplot}
\usepackage{mathtools}
\usetikzlibrary{decorations.pathreplacing}
% Математические пакеты
\usepackage{amsmath, amssymb, amsthm, mathrsfs}
\usepackage{geometry}
\usepackage{graphicx}
\usepackage{pgfplots}
\pgfplotsset{compat=1.18}
\geometry{left=2cm, right=2cm, top=2cm, bottom=2cm}
\usepackage[colorlinks=true, linkcolor=blue, citecolor=black]{hyperref}
\newcommand{\norm}[1]{\left\lVert#1\right\rVert}


\begin{document}
\tableofcontents
\chapter{Билет 1} 
\section{Билет 1. Ортогональные и ортонормированные системы в Гильбертовом пр-ве. Базис. Плотность множеств. Сепарабельность.}
\subsection{Определения}

\noindent \textbf{Опр.} \underline{Евклидово пространство}: $E$ -- линейное пространство со \underline{скалярным произведением}: $(\cdot, \cdot)$ -- билинейная форма ($\mathbb{R}$)\\
\phantom{со скалярным произведением: $(\cdot, \cdot)$ -- } 1,5-линейная форма ($\mathbb{C}$)

\begin{enumerate}
    \item $(\alpha_1 f_1 + \alpha_2 f, g) = \alpha_1(f_1, g) + \alpha_2(f_2, g)$
    \item $(f, g) = (g, f) \quad (\mathbb{R})$ \\
          $(f, g) = \overline{(g, f)} \quad (\mathbb{C})$
    \item $(f, f) \ge 0$ и $(f, f) = 0$ только если $f = 0$.
\end{enumerate}

\bigskip

\noindent Тогда $\norm{f} := \sqrt{(f, f)}$ -- норма

\bigskip

\noindent \underline{Неравенство Коши}: $|(f, g)| \le \norm{f} \cdot \norm{g}$

\bigskip

\noindent \textbf{Д-во:}
$$
\norm{f - \lambda g}^2 = (f - \lambda g, f - \lambda g) = (f, f) - \lambda(g, f) - \overline{\lambda}(f, g) + |\lambda|^2(g, g) =
$$
$$
= \norm{f}^2 - 2\text{Re}(\lambda (g, f)) + |\lambda|^2 \norm{g}^2 \ge 0
$$

\noindent для $\lambda = \dfrac{(f, g)}{\norm{g}^2}$:
$$
\norm{f}^2 - 2 \dfrac{|(f, g)|^2}{\norm{g}^2} + \dfrac{|(f, g)|^2}{\norm{g}^4} \norm{g}^2 =
$$
$$
= \norm{f}^2 - \dfrac{|(f, g)|^2}{\norm{g}^2} \ge 0 \implies \norm{f} \cdot \norm{g} \ge |(f, g)| \quad \square
$$


\textbf{Опр (Гильбертово пр-во):}
Полное евклидово пр-во, снабженное скалярным произведением с размерностью $\infty$.

\textbf{Опр:} $f, g$ ортогональны, если $(f, g) = 0$.
\[
\varphi = \arccos \left( \frac{|(f, g)|}{\|f\| \cdot \|g\|} \right)
\]

\textbf{Опр:} $\{x_\alpha\}_\alpha$ -- ортогональная система, если $\forall \alpha \neq \beta \quad x_\alpha \perp x_\beta$.

\textbf{Замечание:}
Орт $\Rightarrow$ ЛНЗ (линейная независимость).
\[
\left\| \sum x_k \right\|^2 = \sum \|x_k\|^2 \quad (x_\alpha \neq 0)
\]
\[
C_1 x_{\alpha_1} + \dots + C_n x_{\alpha_n} = 0 \Rightarrow \sum |C_k|^2 \|x_{\alpha_k}\|^2 = 0 \Rightarrow C_k = 0, \forall k.
\]

\textbf{Опр (Полная С):}
$\{x_\alpha\}_\alpha$ -- полная, если $\text{span}\{x_\alpha\} = H$.

\textbf{Опр:} $L \subset H$ -- линеал, если $\forall x, y \in L, \forall \alpha, \beta \in \mathbb{C} (\mathbb{R}), \quad \alpha x + \beta y \in L$.

$H_1 \subset H$ -- подпространство, если $H_1$ -- линеал и $\overline{H_1} = H_1$ (замкнутый линеал).

$L$ -- замкнут, если $\forall \{x_n\} \in L$, которая сходится к пределу $x$ ($x_n \to x$), то $x \in L$.

\textbf{Опр:} Полная ортогональная система называется \underline{Базисом} Гильбертова пр-ва.

\textbf{Опр:} $\forall \alpha \ \|x_\alpha\| = 1$ -- ортонормированная система.

\textbf{Опр:} $\{x_\alpha\}$ -- ОНС $\Rightarrow \left\{ \frac{x_\alpha}{\|x_\alpha\|} \right\}$ -- ОНС.
\[
(x_\alpha, x_\beta) = \delta_{\alpha, \beta} = 
\begin{cases} 
1, & \alpha = \beta \\ 
0, & \alpha \neq \beta 
\end{cases}
\]

\textbf{Опр (Плотное множество):}
Множество $A$ -- плотно в топологическом пространстве $X$, если $\overline{A} = X$.

\textbf{Опр:} Если в пр-ве $\exists$ счетное плотное мн-во, пространство -- \underline{сепарабельное}.

\textbf{Теорема:} В сепарабельном пр-ве (евклидовом) любая ортогональная система не более, чем счетна.

\subsection{Док-во:}

Пусть $\varphi_\alpha \perp \varphi_\beta, \forall \alpha \neq \beta$ и $\|\varphi_\alpha\| = 1$ (н.к. -- не ограничиваящ общности / нормированный класс).

\[
\|\varphi_\alpha - \varphi_\beta\|^2 = (\varphi_\alpha - \varphi_\beta, \varphi_\alpha - \varphi_\beta) = (\varphi_\alpha, \varphi_\alpha) - \underbrace{(\varphi_\alpha, \varphi_\beta)}_{=0} - \underbrace{(\varphi_\beta, \varphi_\alpha)}_{=0} + (\varphi_\beta, \varphi_\beta) = ?
\]
Так как ортогональны (смешанные произведения 0) и нормированы (скалярные квадраты 1):
\[
= 1 + 1 = 2.
\]
\[
\Rightarrow \|\varphi_\alpha - \varphi_\beta\| = \sqrt{2}.
\]

Вокруг каждого вектора $\varphi_\alpha$ построим открытый шар радиуса $r < \frac{\sqrt{2}}{2}$.
Т.к. расстояние между центрами шаров $= \sqrt{2}$, они не пересекаются.
\[
B_{\frac{\sqrt{2}}{2}}(\varphi_\alpha) \cap B_{\frac{\sqrt{2}}{2}}(\varphi_\beta) = \emptyset, \quad \forall \alpha \neq \beta.
\]

Т.к. пр-во сепарабельное, в нем есть счетное плотное множество:
\[
A = \{a_1, a_2, \dots\}
\]

По определению плотности, в каждый шар $B(\varphi_\alpha)$ должен попасть хотя бы один элемент из $A$.
\[
\forall \alpha \ \exists a \in A : a \in B(\varphi_\alpha).
\]
$\Rightarrow$ Кол-во шаров не может превосходить количество элементов в $A$.
$A$ -- счетно $\Rightarrow$ система $\{\varphi_\alpha\}$ -- счетная.

\[
[\{\varphi_\alpha\}] \subseteq [A] \quad (\text{счетное множество}).
\]
\chapter{Билет 2 }


\section{Теорема об ортогонализации. Существование базиса в сепарабельном гильбертовом пространстве.}

\textbf{Теорема (об ортогонализации):}
Пусть система $\{f_n\}_{n=1}^\infty$ — линейно независимая (Л.Н.З.) (любой конечный набор векторов из неё Л.Н.З.).

Тогда существует ортонормированная система (О.Н.С.) $\{\varphi_n\}_{n=1}^\infty$, такая что линейные оболочки совпадают на каждом шаге:
\[
L_n = \text{span}\{f_1, \dots, f_n\} = \text{span}\{\varphi_1, \dots, \varphi_n\}.
\]
Векторы новой системы выражаются через старые, и наоборот (матрица перехода треугольная):
\begin{align*}
\varphi_n &= a_{n1} f_1 + \dots + a_{nn} f_n, \quad a_{nn} \neq 0, \\
f_n &= \beta_{n1} \varphi_1 + \dots + \beta_{nn} \varphi_n, \quad \beta_{nn} \neq 0.
\end{align*}

\subsubsection{Доказательство (метод Грама-Шмидта)}

\textbf{Шаг 1 (База индукции):}
Пусть $n=1$. Положим:
\[
\varphi_1 = \frac{f_1}{\|f_1\|}.
\]
Так как исходная система Л.Н.З., то $f_1 \neq 0$, следовательно деление корректно.
Очевидно, что $\text{span}\{f_1\} = \text{span}\{\varphi_1\}$.

\textbf{Шаг 2 (Индуктивный переход):}
Пусть векторы $\varphi_1, \dots, \varphi_{n-1}$ уже построены, они известны, ортонормированы и их линейная оболочка совпадает со $\text{span}\{f_1, \dots, f_{n-1}\}$.
Построим $\varphi_n$.

Будем искать вспомогательный вектор $h_n$ в виде разности вектора $f_n$ и его проекции на уже построенное подпространство (линейной комбинации предыдущих $\varphi_k$):
\[
h_n = f_n - (c_1 \varphi_1 + \dots + c_{n-1} \varphi_{n-1}) = f_n - \sum_{k=1}^{n-1} c_k \varphi_k.
\]
Коэффициенты $c_k$ подбираются так, чтобы $h_n$ был ортогонален всем предыдущим векторам. Умножим скалярно на $\varphi_j$ (где $j < n$):
\[
(h_n, \varphi_j) = (f_n, \varphi_j) - \sum_{k=1}^{n-1} c_k \underbrace{(\varphi_k, \varphi_j)}_{\delta_{kj}} = (f_n, \varphi_j) - c_j \cdot 1.
\]
Требуем $(h_n, \varphi_j) = 0$, отсюда получаем выражение для коэффициентов Фурье:
\[
c_j = (f_n, \varphi_j).
\]
Таким образом, $h_n$ ортогонален всем предыдущим:
\[
(h_n, \varphi_k) = 0, \quad \forall k=1, \dots, n-1.
\]
Так как $\text{span}\{\varphi_j\}_{j=1}^{n-1} = \text{span}\{f_j\}_{j=1}^{n-1}$, то из ортогональности $h_n \perp \varphi_1, \dots, \varphi_{n-1}$ следует:
\[
h_n \perp f_1, \dots, f_{n-1}.
\]
Вектор $h_n \neq 0$, так как иначе $f_n$ выражался бы через $f_1, \dots, f_{n-1}$, что противоречит Л.Н.З. исходной системы.

Нормируем полученный вектор:
\[
\varphi_n = \frac{h_n}{\|h_n\|}.
\]
Из построения видно, что $\varphi_n$ является линейной комбинацией $f_n$ и $\varphi_1, \dots, \varphi_{n-1}$ (а значит и $f_1, \dots, f_{n-1}$).
Выразим обратно $f_n$:
\[
f_n = h_n + \sum_{k=1}^{n-1} c_k \varphi_k = \|h_n\| \varphi_n + \sum_{k=1}^{n-1} c_k \varphi_k.
\]
Здесь коэффициент при старшем члене $\beta_{nn} = \|h_n\| \neq 0$, а $\beta_{nk} = c_k$.
Теорема доказана.

\subsection{Следствие (Существование базиса)}

\textbf{Утверждение:}
Если пространство $E$ — сепарабельно, то в $E$ существует ортонормированный базис (О.Н.Б.).

\textbf{Доказательство:}
1. Раз $E$ — сепарабельно, то по определению в нем существует счетное плотное множество $A$:
\[
A = \{a_n\}_{n=1}^\infty, \quad \overline{A} = E \quad (\text{где } A = \{a_n \mid n \in \mathbb{N}\}).
\]
2. Система $A$ может быть линейно зависимой (Л.З.).
Идем по порядку и выкидываем линейно зависимые векторы (те, что выражаются через предыдущие).
Получим подсистему $\{\tilde{a}_n\}$, которая является линейно независимой (Л.Н.З.).
При этом линейная оболочка не уменьшилась:
\[
\overline{\text{span}\{\tilde{a}_n\}} = E.
\]
3. Применяем к системе $\{\tilde{a}_n\}$ процедуру ортогонализации Грама-Шмидта.
Получаем ортонормированную систему (О.М.С.) $\{\varphi_n\}$.
Так как $\text{span}\{\varphi_n\} = \text{span}\{\tilde{a}_n\}$, то замыкание линейной оболочки $\{\varphi_n\}$ совпадает с $E$:
\[
\overline{\text{span}\{\varphi_n\}} = E.
\]
По определению, полная ортонормированная система является \textbf{ортонормированным базисом}.
\chapter{Билет 3}



\section{Ряд Фурье в гильбертовом пространстве. Неравенство Бесселя. Замкнутость.}

Пусть $\{\varphi_n\}_{n=1}^\infty$ — ортонормированная система (О.Н.С.) в евклидовом пространстве $E$.
Пусть $f \in E$.

Числа $c_n = (f, \varphi_n)$ называются \textbf{коэффициентами Фурье} элемента $f$ по системе $\{\varphi_n\}$.

Формальный ряд $\sum_{n=1}^\infty c_n \varphi_n$ называется \textbf{рядом Фурье} элемента $f$ по системе $\{\varphi_n\}$.

Обозначим через $S_n$ частичную сумму ряда Фурье:
\[
S_n = \sum_{k=1}^n c_k \varphi_k.
\]

\subsection{Оценка квадрата нормы разности}

Оценим квадрат нормы разности между функцией и частичной суммой: $\|f - S_n\|^2$.
Используя свойства скалярного произведения (линейность и сопряженную симметрию), раскроем скобки:
\[
\|f - S_n\|^2 = (f - S_n, f - S_n) = \|f\|^2 - (f, S_n) - (S_n, f) + \|S_n\|^2.
\]

Рассмотрим слагаемые отдельно:
\begin{enumerate}
    \item Скалярное произведение $(f, S_n)$:
    \[
    (f, S_n) = \left( f, \sum_{k=1}^n c_k \varphi_k \right) = \sum_{k=1}^n \overline{c_k} (f, \varphi_k) = \sum_{k=1}^n \overline{c_k} c_k = \sum_{k=1}^n |c_k|^2.
    \]
    
    \item Скалярное произведение $(S_n, f)$:
    \[
    (S_n, f) = \overline{(f, S_n)} = \sum_{k=1}^n c_k \overline{c_k} = \sum_{k=1}^n |c_k|^2.
    \]
    
    \item Квадрат нормы $\|S_n\|^2$:
    \[
    \|S_n\|^2 = \left( \sum_{k=1}^n c_k \varphi_k, \sum_{k=1}^n c_k \varphi_k \right).
    \]
    Так как система ортонормирована ($\varphi_i \perp \varphi_j$ при $i \neq j$ и $\|\varphi_i\|=1$), работает обобщенная теорема Пифагора:
    \[
    \|S_n\|^2 = \sum_{k=1}^n |c_k|^2.
    \]
\end{enumerate}

Подставим полученные значения обратно в исходное равенство:
\[
\|f - S_n\|^2 = \|f\|^2 - \sum_{k=1}^n |c_k|^2 - \sum_{k=1}^n |c_k|^2 + \sum_{k=1}^n |c_k|^2.
\]
В итоге получаем важное тождество:
\[
\|f - S_n\|^2 = \|f\|^2 - \sum_{k=1}^n |c_k|^2.
\]

\subsection{Неравенство Бесселя}

Так как квадрат нормы вектора всегда неотрицателен ($\|f - S_n\|^2 \ge 0$), то:
\[
\|f\|^2 - \sum_{k=1}^n |c_k|^2 \ge 0 \implies \sum_{k=1}^n |c_k|^2 \le \|f\|^2.
\]
Это неравенство справедливо для любого $n$. Перейдем к пределу при $n \to \infty$. Так как частичные суммы ограничены сверху числом $\|f\|^2$ и монотонно возрастают (сумма неотрицательных чисел), ряд сходится:
\[
\sum_{n=1}^\infty |c_n|^2 \le \|f\|^2.
\]
Это неравенство называется \textbf{неравенством Бесселя}.

\textbf{Следствие:}
Так как ряд $\sum |c_n|^2$ сходится, то его общий член стремится к нулю:
\[
|c_n| \to 0 \quad \text{при } n \to \infty.
\]

\subsection{Геометрический смысл (Минимальное свойство)}

Вопрос: Как наилучшим образом приблизить вектор $f$ линейной комбинацией векторов $\varphi_1, \dots, \varphi_n$?

Частичная сумма Фурье $S_n = \sum_{k=1}^n c_k \varphi_k$ является \textbf{ортогональной проекцией} вектора $f$ на линейную оболочку (подпространство) $\text{span}\{\varphi_1, \dots, \varphi_n\}$.

Вектор ошибки ортогонален подпространству:
\[
\left( f - \sum_{k=1}^n c_k \varphi_k \right) \perp \text{span}\{\varphi_1, \dots, \varphi_n\}.
\]

\begin{center}
\begin{tikzpicture}[scale=2, >=latex]
    % Axes
    \draw[->] (0,0,0) -- (1.5,0,0) node[anchor=north east]{$\varphi_2$};
    \draw[->] (0,0,0) -- (0,1.5,0) node[anchor=north west]{$\varphi_3$};
    \draw[->] (0,0,0) -- (0,0,1.5) node[anchor=south]{$\varphi_1$};
    
    % Vector f
    \draw[thick, ->, blue] (0,0,0) -- (1.2, 1.2, 0.5) node[right]{$f$};
    
    % Plane representing span
    \fill[gray!10, opacity=0.5] (0,0,0) -- (1.5,0,0) -- (1.5,0,1.5) -- (0,0,1.5) -- cycle;
    
    % Projection Sn
    \draw[thick, ->, red] (0,0,0) -- (1.2, 0, 0.5) node[below right]{$S_n = \sum c_k \varphi_k$};
    
    % Error vector
    \draw[dashed, ->] (1.2, 0, 0.5) -- (1.2, 1.2, 0.5) node[midway, right]{$f - S_n$};
    
    % Right angle symbol roughly
    \draw (1.2, 0, 0.5) -- ++(0, 0.1, 0) -- ++(-0.1, 0, 0);

\end{tikzpicture}
\end{center}

\section{Замкнутость системы}

\textbf{Определение:}
Ортонормированная система (О.Н.С.) называется \textbf{замкнутой}, если для любого $f \in E$ неравенство Бесселя обращается в равенство:
\[
\sum_{k=1}^\infty |c_k|^2 = \|f\|^2.
\]
Это равенство называется \textbf{равенством Парсеваля}.

\textbf{Следствие:}
Если система замкнута, то:
\[
\|f - S_n\|^2 = \|f\|^2 - \sum_{k=1}^n |c_k|^2 \to 0 \quad (\text{при } n \to \infty).
\]
Это означает, что ряд Фурье сходится к самой функции $f$ по норме пространства $E$:
\[
\sum_{n=1}^\infty c_n \varphi_n = f.
\]
\chapter{Билет 4 }
\section{Равносильность замкнутости и полноты ортогональной системы в гильбертовом пространстве.}

\textbf{Определение (Замкнутость):}
Ортонормированная система (О.Н.С.) $\{\varphi_n\}$ называется \textbf{замкнутой}, если для любого $f$ из пространства выполняется равенство Парсеваля:
\[
\sum_{n=1}^\infty |c_n|^2 = \|f\|^2, \quad \text{где } c_n = (f, \varphi_n).
\]

\textbf{Определение (Полнота):}
О.Н.С. $\{\varphi_n\}$ называется \textbf{полной}, если в пространстве $E$ не существует ненулевого вектора, ортогонального всем векторам системы.
\[
(\forall n \ (f, \varphi_n) = 0) \implies f = 0.
\]

\textbf{Теорема (Система замкнута $\iff$ полна):}
В сепарабельном гильбертовом пространстве для ортонормированной системы свойство быть замкнутой равносильно свойству быть полной.

\subsection{Доказательство}

\textbf{1. Замкнутость $\implies$ Полнота ($\Rightarrow$)}

Пусть система $\{\varphi_n\}$ — замкнута. Это означает, что для любого $f \in E$ ряд Фурье сходится к функции по норме ($S_n \to f$) и выполнено равенство Парсеваля.

Предположим, что существует вектор $f \in E$, который ортогонален всей системе:
\[
f \perp \{\varphi_n\} \quad \forall n.
\]
Вычислим его коэффициенты Фурье:
\[
c_n = (f, \varphi_n) = 0 \quad \forall n.
\]
Так как система замкнута, для вектора $f$ верно равенство Парсеваля:
\[
\|f\|^2 = \sum_{n=1}^\infty |c_n|^2 = \sum_{n=1}^\infty 0 = 0.
\]
Если квадрат нормы равен 0, то и сам вектор равен 0:
\[
\|f\|^2 = 0 \implies f = 0.
\]
\textbf{Вывод:} Единственный вектор, перпендикулярный системе — это ноль. Следовательно, система — полна.

\vspace{1em}

\textbf{2. Полнота $\implies$ Замкнутость ($\Leftarrow$)}

Так как система полна, то линейная оболочка системы плотна в $E$. Это значит, что для любого $\varepsilon > 0$ и любого $f \in E$ существуют $n$ и коэффициенты $a_1, \dots, a_n$, такие, что линейная комбинация $\Phi_n = \sum_{k=1}^n a_k \varphi_k$ приближает вектор $f$ с точностью $\varepsilon$:
\[
\left\| f - \sum_{k=1}^n a_k \varphi_k \right\| < \varepsilon.
\]

Сравним это приближение с частичной суммой ряда Фурье $S_n$.
Рассмотрим квадрат нормы разности:
\[
\left\| f - \sum_{k=1}^n a_k \varphi_k \right\|^2 = \left( f - \sum_{k=1}^n a_k \varphi_k, f - \sum_{l=1}^n a_l \varphi_l \right).
\]
Раскроем скобки, используя линейность скалярного произведения:
\[
= (f, f) - \left( f, \sum_{l=1}^n a_l \varphi_l \right) - \left( \sum_{k=1}^n a_k \varphi_k, f \right) + \left( \sum_{k=1}^n a_k \varphi_k, \sum_{l=1}^n a_l \varphi_l \right).
\]
Вычислим каждое слагаемое:
\begin{enumerate}
    \item $(f, f) = \|f\|^2$.
    \item $\left( f, \sum_{l=1}^n a_l \varphi_l \right) = \sum_{l=1}^n \overline{a_l} (f, \varphi_l) = \sum_{l=1}^n \overline{a_l} c_l$ (так как $c_l = (f, \varphi_l)$).
    \item $\left( \sum_{k=1}^n a_k \varphi_k, f \right) = \sum_{k=1}^n a_k ( \varphi_k, f ) = \sum_{k=1}^n a_k \overline{c_k}$.
    \item Для последнего слагаемого используем ортонормированность $(\varphi_k, \varphi_l) = \delta_{kl} = \begin{cases} 1, & k=l \\ 0, & k \neq l \end{cases}$:
    \[
    \left( \sum_{k=1}^n a_k \varphi_k, \sum_{l=1}^n a_l \varphi_l \right) = \sum_{k=1}^n \sum_{l=1}^n a_k \overline{a_l} (\varphi_k, \varphi_l) = \sum_{k=1}^n a_k \overline{a_k} \cdot 1 = \sum_{k=1}^n |a_k|^2.
    \]
\end{enumerate}

Подставим всё обратно:
\[
\| f - \Phi_n \|^2 = \|f\|^2 - \sum_{k=1}^n \overline{a_k} c_k - \sum_{k=1}^n a_k \overline{c_k} + \sum_{k=1}^n |a_k|^2.
\]

Добавим и вычтем сумму квадратов модулей коэффициентов Фурье $\sum_{k=1}^n |c_k|^2$:
\[
= \left( \|f\|^2 - \sum_{k=1}^n |c_k|^2 \right) + \left( \sum_{k=1}^n |c_k|^2 - \sum_{k=1}^n \overline{a_k} c_k - \sum_{k=1}^n a_k \overline{c_k} + \sum_{k=1}^n |a_k|^2 \right).
\]
Первая скобка — это $\|f - S_n\|^2$. Вторую свернем в квадрат разности:
\[
= \|f - S_n\|^2 + \sum_{k=1}^n (c_k - a_k)(\overline{c_k} - \overline{a_k}) = \|f - S_n\|^2 + \sum_{k=1}^n |c_k - a_k|^2.
\]
Получили тождество:
\[
\left\| f - \sum_{k=1}^n a_k \varphi_k \right\|^2 = \|f - S_n\|^2 + \underbrace{\sum_{k=1}^n |c_k - a_k|^2}_{\ge 0}.
\]
Отсюда следует неравенство (экстремальное свойство коэффициентов Фурье):
\[
\|f - S_n\|^2 \le \left\| f - \sum_{k=1}^n a_k \varphi_k \right\|^2.
\]
По условию полноты существует такой набор $a_k$, что ошибка аппроксимации меньше $\varepsilon$:
\[
\left\| f - \sum_{k=1}^n a_k \varphi_k \right\| < \varepsilon.
\]
Из неравенства выше следует:
\[
\|f - S_n\| < \varepsilon.
\]
Так как неравенство верно для любого $\varepsilon > 0$, перейдем к пределу при $n \to \infty$:
\[
\lim_{n \to \infty} \|f - S_n\|^2 = 0.
\]
Или, подставляя выражение для нормы разности:
\[
\lim_{n \to \infty} \left( \|f\|^2 - \sum_{k=1}^n |c_k|^2 \right) = 0 \implies \|f\|^2 = \sum_{k=1}^\infty |c_k|^2.
\]
Это и есть равенство Парсеваля.
\textbf{Следовательно, система замкнута.}

\chapter{Билет 5}
\section{ Теорема Рисса — Фишера. Критерий полноты ортогональной системы в сепарабельном гильбертовом простран-
стве.}
\subsection{Теорема (Рисса — Фишера)}

\textbf{Условие:}
Пусть $H$ — сепарабельное гильбертово пространство, и $\{\varphi_n\}_{n=1}^\infty$ — ортонормированная система (О.Н.С.) в $H$.

Для любой последовательности чисел $\{c_k\}_{k=1}^\infty$, принадлежащей пространству $l^2$ (то есть последовательности, для которой сходится ряд квадратов модулей: $\sum_{k=1}^\infty |c_k|^2 < \infty$), существует элемент $f \in H$, такой что:
\begin{enumerate}
    \item Числа $c_k$ являются его коэффициентами Фурье: $c_k = (f, \varphi_k)$.
    \item Выполняется равенство Парсеваля: $\|f\|^2 = \sum_{k=1}^\infty |c_k|^2$.
\end{enumerate}

\textbf{Доказательство:}

Рассмотрим последовательность частичных сумм ряда:
\[
f_n = \sum_{k=1}^n c_k \varphi_k.
\]
Докажем, что эта последовательность сходится к некоторому вектору $f \in H$.

\textbf{1. Проверка фундаментальности (Критерий Коши).}
Рассмотрим квадрат нормы разности двух частичных сумм при $m > n$:
\[
\|f_m - f_n\|^2 = \left\| \sum_{k=n+1}^m c_k \varphi_k \right\|^2 = \left( \sum_{k=n+1}^m c_k \varphi_k, \sum_{l=n+1}^m c_l \varphi_l \right).
\]
Раскрывая скалярное произведение и используя ортонормированность системы $(\varphi_k, \varphi_l) = \delta_{kl}$:
\[
= \sum_{k=n+1}^m \sum_{l=n+1}^m c_k \overline{c_l} (\varphi_k, \varphi_l) = \sum_{k=n+1}^m |c_k|^2.
\]
Так как исходный ряд $\sum |c_k|^2$ сходится (по условию $\{c_k\} \in l^2$), то для любого $\varepsilon > 0$ существует номер $N \in \mathbb{N}$, такой что для всех $m, n > N$:
\[
\sum_{k=n+1}^m |c_k|^2 < \varepsilon.
\]
Следовательно, последовательность $\{f_n\}$ — фундаментальная.

\textbf{2. Существование предела.}
Так как $H$ — гильбертово пространство, оно является полным. Значит, любая фундаментальная последовательность в нем имеет предел:
\[
\exists \lim_{n \to \infty} f_n = f \in H.
\]

\textbf{3. Проверка коэффициентов Фурье.}
Докажем, что $c_k$ действительно являются коэффициентами Фурье вектора $f$. Рассмотрим скалярное произведение $(f, \varphi_k)$. В силу непрерывности скалярного произведения:
\[
(f, \varphi_k) = (\lim_{n \to \infty} f_n, \varphi_k) = \lim_{n \to \infty} (f_n, \varphi_k).
\]
Рассмотрим $(f_n, \varphi_k)$ при $n > k$:
\[
f_n = c_1 \varphi_1 + \dots + c_k \varphi_k + \dots + c_n \varphi_n.
\]
Из ортогональности $(f_n, \varphi_k) = c_k$ (все остальные слагаемые зануляются).
Переходя к пределу при $n \to \infty$, получаем:
\[
(f, \varphi_k) = c_k.
\]

\textbf{4. Равенство Парсеваля.}
\[
\|f\|^2 = \lim_{n \to \infty} \|f_n\|^2 = \lim_{n \to \infty} \sum_{k=1}^n |c_k|^2 = \sum_{k=1}^\infty |c_k|^2.
\]
Теорема доказана.

\subsubsection{Следствие (Изоморфизм пространств):}
Сепарабельное гильбертово пространство $H$ изометрически изоморфно пространству последовательностей $l^2$:
\[
f \in H \cong \{c_k\} \in l^2.
\]

\subsection{Теорема (Критерий полноты)}

\textbf{Условие:}
Пусть $H$ — сепарабельное гильбертово пространство. Ортонормированная система $\{\varphi_n\}$ является \textbf{полной} тогда и только тогда, когда в $H$ не существует ненулевого вектора, ортогонального всем векторам системы.
\[
\{\varphi_n\} \text{ — полна } \iff \left( \forall g \in H: (\forall n \ (g, \varphi_n) = 0) \implies g = 0 \right).
\]

\textbf{Доказательство:}

\textbf{1. Необходимость ($\Rightarrow$)}
Пусть система полна. Тогда (по предыдущей теореме) она является замкнутой.
Предположим, существует $g \in H$, такой что $(g, \varphi_n) = 0$ для всех $n$.
Тогда все коэффициенты Фурье $c_n(g) = 0$.
Так как система замкнута, верно равенство Парсеваля:
\[
\|g\|^2 = \sum_{n=1}^\infty |c_n|^2 = \sum 0 = 0.
\]
Следовательно, $g = 0$.

\textbf{2. Достаточность ($\Leftarrow$)}
Докажем от противного.
Пусть условие отсутствия ортогонального вектора выполнено, но система \textbf{не полна}.
Тогда она и не замкнута (так как в гильбертовом пр-ве полнота $\Leftrightarrow$ замкнутость).
Это означает, что существует некоторый вектор $h \in H$, для которого нарушается равенство Парсеваля (верно строгое неравенство Бесселя):
\[
\sum_{n=1}^\infty |c_n|^2 < \|h\|^2, \quad \text{где } c_n = (h, \varphi_n).
\]
Возьмем последовательность коэффициентов $\{c_n\}$. Она принадлежит $l^2$.
По теореме Рисса — Фишера существует вектор $f \in H$, такой что:
\[
c_n(f) = c_n \quad \text{и} \quad \|f\|^2 = \sum_{n=1}^\infty |c_n|^2.
\]
Сравним векторы $f$ и $h$. Рассмотрим разность $g = f - h$.
\begin{enumerate}
    \item Вектор $g$ не равен нулю, так как нормы $f$ и $h$ различны ($\|f\| < \|h\|$).
    \item Проверим ортогональность вектора $g$ системе $\{\varphi_n\}$:
    \[
    (g, \varphi_n) = (f - h, \varphi_n) = (f, \varphi_n) - (h, \varphi_n) = c_n - c_n = 0.
    \]
\end{enumerate}
Мы нашли ненулевой вектор $g \neq 0$, который ортогонален всем векторам системы.
Это противоречит условию теоремы (что таких векторов не существует).

Следовательно, предположение неверно, и система полна.

\chapter{Билет 6}
\section{Пространство \text{$L_2([-\pi, \pi])$} и ортогональная система}

Рассмотрим пространство функций, интегрируемых с квадратом на отрезке $[-\pi, \pi]$, обозначаемое как $L_2([-\pi, \pi])$. Функции в этом пространстве интегрируемы по Лебегу.

Скалярное произведение двух функций $f$ и $g$ в этом пространстве определяется следующим образом:
$$
(f, g) = \int_{-\pi}^{\pi} f(x) \overline{g(x)} \, dx
$$
(в вещественном случае сопряжение можно опустить: $f(x)g(x)$).

Норма функции определяется через скалярное произведение:
$$
\|f\| = \sqrt{(f, f)} = \sqrt{\int_{-\pi}^{\pi} |f(x)|^2 \, dx}
$$

В этом пространстве рассмотрим тригонометрическую систему функций:
$$
\{ 1, \sin x, \cos x, \sin 2x, \cos 2x, \dots, \sin nx, \cos nx, \dots \}
$$

\subsection{Доказательство ортогональности системы}

Эта система является ортогональной. Докажем это, вычислив скалярные произведения различных элементов базиса.

1.  **Интеграл от произведения синуса и косинуса:**
    $$
    \int_{-\pi}^{\pi} \sin nx \cos mx \, dx = 0
    $$
    Это равенство верно для любых $n$ и $m$, так как подынтегральная функция является нечетной (произведение нечетной функции $\sin$ и четной $\cos$ есть нечетная функция), а интегрирование ведется по симметричному промежутку.

2.  **Интеграл от произведения косинусов ($n \neq m$):**
    Воспользуемся формулой произведения косинусов:
    $$
    \cos nx \cos mx = \frac{1}{2} [\cos((n-m)x) + \cos((n+m)x)]
    $$
    Тогда интеграл:
    $$
    \int_{-\pi}^{\pi} \cos nx \cos mx \, dx = \frac{1}{2} \int_{-\pi}^{\pi} \cos((n-m)x) \, dx + \frac{1}{2} \int_{-\pi}^{\pi} \cos((n+m)x) \, dx
    $$
    Так как интеграл от косинуса по полному периоду (или кратному ему) равен нулю:
    $$
    \left. \frac{\sin((n-m)x)}{n-m} \right|_{-\pi}^{\pi} = 0 \quad \text{и} \quad \left. \frac{\sin((n+m)x)}{n+m} \right|_{-\pi}^{\pi} = 0
    $$
    Следовательно, при $n \neq m$ скалярное произведение равно 0.

3.  **Вычисление квадратов норм (нормировка):**
    При $n = m$ имеем:
    $$
    \int_{-\pi}^{\pi} \sin^2(nx) \, dx = \int_{-\pi}^{\pi} \cos^2(nx) \, dx = \pi
    $$
    Это следует из того, что $\sin^2 \alpha + \cos^2 \alpha = 1$, и интегралы от них равны. Сумма интегралов равна длине интервала $2\pi$, значит каждый равен $\pi$.
    
    Для единицы (свободного члена):
    $$
    \int_{-\pi}^{\pi} 1^2 \, dx = 2\pi \Rightarrow \|1\|^2 = 2\pi.
    $$

Чтобы получить **ортонормированный базис (ОНБ)**, разделим функции на их нормы:
$$
\frac{1}{\sqrt{2\pi}}, \quad \frac{\cos nx}{\sqrt{\pi}}, \quad \frac{\sin nx}{\sqrt{\pi}} \quad (n \ge 1)
$$

\section{Ряд Фурье в вещественной форме}

Если функция $f \in L_2([-\pi, \pi])$, то мы можем сопоставить ей ряд Фурье:
$$
f(x) \sim \frac{a_0}{2} + \sum_{n=1}^{\infty} (a_n \cos nx + b_n \sin nx)
$$
Коэффициенты ряда определяются формулами Фурье:
$$
a_n = \frac{1}{\pi} \int_{-\pi}^{\pi} f(x) \cos nx \, dx \quad (n=0, 1, 2, \dots)
$$
$$
b_n = \frac{1}{\pi} \int_{-\pi}^{\pi} f(x) \sin nx \, dx \quad (n=1, 2, \dots)
$$
\textit{Примечание:} Коэффициент $a_0$ вычисляется по общей формуле для $a_n$, но в разложении делится на 2, чтобы компенсировать отличие нормы функции $1$ (которая равна $\sqrt{2\pi}$) от нормы $\cos nx$ (которая равна $\sqrt{\pi}$).

\textbf{Сходимость:}
Для любой функции $f \in L_2([-\pi, \pi])$ частичные суммы ряда Фурье $S_n(x)$ сходятся к функции $f(x)$ по норме пространства $L_2$ (сходимость в среднем квадратичном):
$$
\| f - S_n \| \to 0 \quad \text{при} \quad n \to \infty.
$$
Также стоит отметить, что $L_2([-\pi, \pi]) \subset L_1([-\pi, \pi])$, поэтому коэффициенты определены корректно.

\section{Комплексная форма ряда Фурье}

Используя формулу Эйлера $e^{inx} = \cos nx + i \sin nx$, можно перейти к комплексной форме записи.

Рассмотрим систему функций:
$$
\{ e^{inx} \}_{n \in \mathbb{Z}} = \{ \dots, e^{-ix}, 1, e^{ix}, \dots \}
$$
Это полная ортогональная система в $L_2$.

\textbf{Ортогональность:}
Проверим скалярное произведение (учитывая комплексное сопряжение во втором множителе):
$$
(e^{inx}, e^{imx}) = \int_{-\pi}^{\pi} e^{inx} \overline{e^{imx}} \, dx = \int_{-\pi}^{\pi} e^{inx} e^{-imx} \, dx = \int_{-\pi}^{\pi} e^{i(n-m)x} \, dx
$$
Если $n \neq m$:
$$
= \left. \frac{e^{i(n-m)x}}{i(n-m)} \right|_{-\pi}^{\pi} = \frac{(-1)^{n-m} - (-1)^{n-m}}{i(n-m)} = 0
$$
Если $n = m$:
$$
\|e^{inx}\|^2 = \int_{-\pi}^{\pi} e^{i \cdot 0} \, dx = \int_{-\pi}^{\pi} 1 \, dx = 2\pi
$$

\textbf{Разложение и коэффициенты:}
Ряд имеет вид:
$$
f(x) = \sum_{n=-\infty}^{+\infty} c_n e^{inx}
$$
Коэффициенты $c_n$ вычисляются как проекция на базис:
$$
c_n = \frac{(f, e^{inx})}{\|e^{inx}\|^2} = \frac{1}{2\pi} \int_{-\pi}^{\pi} f(x) e^{-inx} \, dx
$$

\textbf{Равенство Парсеваля для комплексной формы:}
$$
\frac{1}{2\pi} \int_{-\pi}^{\pi} |f(x)|^2 \, dx = \sum_{n=-\infty}^{+\infty} |c_n|^2
$$

\section{Разложение на интервале $[0, \pi]$}

Пусть функция $f(x)$ задана только на интервале $[0, \pi]$. Мы можем разложить её в ряд Фурье, доопределив её на интервал $[-\pi, 0]$ двумя основными способами: четным или нечетным.

\subsection{Разложение по косинусам (четное продолжение)}

Рассмотрим функцию $\tilde{f}(x)$, доопределенную четным образом:
$$
\tilde{f}(x) = \begin{cases} 
f(x), & x \in [0, \pi] \\
f(-x), & x \in [-\pi, 0]
\end{cases}
$$
Очевидно, что $\tilde{f} \in L_2([-\pi, \pi])$.

\begin{center}
\begin{tikzpicture}[scale=0.8]
    % Axes
    \draw[->] (-4,0) -- (4,0) node[right] {$x$};
    \draw[->] (0,-0.5) -- (0,3) node[above] {$f(x)$};
    
    % Ticks
    \draw[thick] (3, 0.1) -- (3, -0.1) node[below] {$\pi$};
    \draw[thick] (-3, 0.1) -- (-3, -0.1) node[below] {$-\pi$};
    \node at (-0.2, -0.3) {$0$};

    % Vertical bounds
    \draw[dashed] (3, 0) -- (3, 2.5);
    \draw[dashed] (-3, 0) -- (-3, 2.5);

    % Function right side (original)
    \draw[thick, blue] (0, 1) to[out=10, in=180] (1.5, 0.5) to[out=0, in=200] (3, 2.5);
    
    % Function left side (even extension)
    \draw[thick, blue] (0, 1) to[out=170, in=0] (-1.5, 0.5) to[out=180, in=-20] (-3, 2.5);

    \node[blue] at (1.5, 2.5) {Исходная $f(x)$};
    \node[blue] at (-1.5, 2.5) {Продолжение};
\end{tikzpicture}
\end{center}

Так как функция $\tilde{f}(x)$ — четная, ряд Фурье не содержит нечетных функций (синусов), то есть $b_n = 0$.
Действительно:
$$
b_n = \frac{1}{\pi} \int_{-\pi}^{\pi} \underbrace{\tilde{f}(x)}_{\text{четн.}} \underbrace{\sin nx}_{\text{нечетн.}} \, dx = 0 \quad (\text{интеграл от нечетной функции})
$$

Коэффициенты $a_n$ вычисляются следующим образом:
$$
a_n = \frac{1}{\pi} \int_{-\pi}^{\pi} \tilde{f}(x) \cos nx \, dx = \frac{2}{\pi} \int_{0}^{\pi} f(x) \cos nx \, dx
$$
Ряд имеет вид:
$$
f(x) = \frac{a_0}{2} + \sum_{n=1}^{\infty} a_n \cos nx, \quad x \in [0, \pi]
$$

\subsection{Разложение по синусам (нечетное продолжение)}
Аналогично, если доопределить функцию нечетным образом ($f(-x) = -f(x)$), то $a_n = 0$, и ряд будет содержать только синусы:
$$
f(x) = \sum_{n=1}^{\infty} b_n \sin nx, \quad b_n = \frac{2}{\pi} \int_{0}^{\pi} f(x) \sin nx \, dx
$$

\section{Обобщение на интервал $[-l, l]$}

Пусть функция $f(x)$ задана на произвольном симметричном отрезке $[-l, l]$.
Сделаем замену переменной для сведения к интервалу $[-\pi, \pi]$:
$$
t = \frac{\pi x}{l} \implies x = \frac{l t}{\pi}, \quad dx = \frac{l}{\pi} dt
$$
При $x = \pm l$ получаем $t = \pm \pi$.

Тригонометрическая система примет вид:
$$
1, \quad \cos \frac{\pi n x}{l}, \quad \sin \frac{\pi n x}{l}
$$

Ряд Фурье на $[-l, l]$:
$$
f(x) = \frac{a_0}{2} + \sum_{n=1}^{\infty} \left( a_n \cos \frac{\pi n x}{l} + b_n \sin \frac{\pi n x}{l} \right)
$$

Коэффициенты вычисляются по формулам:
$$
a_n = \frac{1}{l} \int_{-l}^{l} f(x) \cos \frac{\pi n x}{l} \, dx
$$
$$
b_n = \frac{1}{l} \int_{-l}^{l} f(x) \sin \frac{\pi n x}{l} \, dx
$$
\chapter{Билет 7}
\section{Лемма Римана}

\section{Формулировка:}
Пусть функция $f(x) \in L_1([a, b])$ (абсолютно интегрируема).
Тогда:
$$ \int_a^b f(x) \sin(\omega x + \alpha) \, dx \xrightarrow[\omega \to \infty]{} 0 $$
(Аналогично для $\cos(\omega x + \alpha)$).

\vspace{0.5cm}

\textbf{Геометрическая интерпретация:}
При больших $\omega$ функция $\sin(\omega x + \alpha)$ очень быстро осциллирует. Если $f(x)$ меняется медленно по сравнению с синусом, то положительные и отрицательные площади под графиком произведения $f(x) \sin(\omega x)$ взаимно уничтожаются, и интеграл стремится к нулю.

\begin{center}
\begin{tikzpicture}
\begin{axis}[
    width=10cm, height=4cm,
    axis lines=middle,
    xtick={0, 6.28}, xticklabels={$a$, $b$},
    ytick=\empty,
    xlabel=$x$, ylabel={},
    domain=0:6.28, samples=200,
    title={Быстрая осцилляция $\sin(\omega x)$ модулируемая $f(x)$}
]
    % Медленная огибающая f(x)
    \addplot[dashed, thick, blue] {exp(-0.2*x)};
    \addplot[dashed, thick, blue] {-exp(-0.2*x)};
    % Произведение f(x) * sin(wx)
    \addplot[red, thick] {exp(-0.2*x) * sin(deg(10*x))};
\end{axis}
\end{tikzpicture}
\end{center}

\section{Доказательство}

Доказательство проводится в два этапа.

\subsection*{1. Пусть $f \in C^1([a, b])$}
Это означает, что функция $f(x)$ и ее производная $f'(x)$ непрерывны на отрезке $[a, b]$.

Рассмотрим интеграл:
$$ I(\omega) = \int_a^b f(x) \sin(\omega x + \alpha) \, dx $$

Воспользуемся методом интегрирования по частям. Введем обозначения:
\begin{align*}
    u &= f(x) \quad \implies \quad du = f'(x) \, dx \\
    dv &= \sin(\omega x + \alpha) \, dx \quad \implies \quad v = -\frac{1}{\omega} \cos(\omega x + \alpha)
\end{align*}
Применяем формулу $\int_a^b u \, dv = uv \big|_a^b - \int_a^b v \, du$:
$$ \int_a^b f(x) \sin(\omega x + \alpha) \, dx = \left[ -\frac{f(x)}{\omega} \cos(\omega x + \alpha) \right]_a^b - \int_a^b \left( -\frac{1}{\omega} \cos(\omega x + \alpha) \right) f'(x) \, dx $$

Раскроем скобки и вынесем константы:
$$ = \frac{f(a)\cos(\omega a + \alpha) - f(b)\cos(\omega b + \alpha)}{\omega} + \frac{1}{\omega} \int_a^b f'(x) \cos(\omega x + \alpha) \, dx \quad (**)$$

\textbf{Оценки:}
Так как $f \in C^1([a, b])$, то функции $f(x)$ и $f'(x)$ ограничены на замкнутом отрезке $[a, b]$. Следовательно, существуют константы $M_1, M_2 > 0$ такие, что:
$$ |f(x)| \le M_1 \quad \forall x \in [a, b] $$
$$ |f'(x)| \le M_2 \quad \forall x \in [a, b] $$

Оценим слагаемые в выражении $(**)$ по модулю:

\textbf{А) Внеинтегральное слагаемое:}
$$ \left| \frac{f(a)\cos(\dots) - f(b)\cos(\dots)}{\omega} \right| \le \frac{|f(a)| \cdot 1 + |f(b)| \cdot 1}{\omega} \le \frac{M_1 + M_1}{\omega} = \frac{2M_1}{\omega} $$

\textbf{Б) Интегральное слагаемое:}
$$ \left| \frac{1}{\omega} \int_a^b f'(x) \cos(\omega x + \alpha) \, dx \right| \le \frac{1}{\omega} \int_a^b |f'(x)| \cdot |\cos(\omega x + \alpha)| \, dx $$
Учитывая, что $|\cos(\dots)| \le 1$ и $|f'(x)| \le M_2$:
$$ \le \frac{1}{\omega} \int_a^b M_2 \cdot 1 \, dx = \frac{M_2}{\omega} (b - a) $$

\textbf{Итоговая оценка:}
$$ \left| \int_a^b f(x) \sin(\omega x + \alpha) \, dx \right| \le \frac{2M_1}{\omega} + \frac{M_2(b - a)}{\omega} $$
При $\omega \to \infty$ правая часть стремится к 0.
$$ \frac{2M_1 + M_2(b-a)}{\omega} \xrightarrow[\omega \to \infty]{} 0 $$
Следовательно, и сам интеграл стремится к нулю.

\vspace{1cm}

\subsection*{2. Пусть $f \in L_1([a, b])$ (Общий случай)}

Используем свойство плотности непрерывно дифференцируемых функций в пространстве $L_1$.
Для любого $\varepsilon > 0$ существует функция $g \in C^1([a, b])$ (гладкая аппроксимация) такая, что норма разности в $L_1$ мала:
$$ \| f - g \|_{L_1([a, b])} = \int_a^b |f(x) - g(x)| \, dx < \frac{\varepsilon}{2} $$

Рассмотрим исходный интеграл, добавив и вычтя $g(x)$:
$$ \int_a^b f(x) \sin(\omega x + \alpha) \, dx = \int_a^b \big( f(x) - g(x) + g(x) \big) \sin(\omega x + \alpha) \, dx $$
Разобьем на два интеграла:
$$ = \underbrace{\int_a^b (f(x) - g(x)) \sin(\omega x + \alpha) \, dx}_{I_1} + \underbrace{\int_a^b g(x) \sin(\omega x + \alpha) \, dx}_{I_2} $$

Оценим каждое слагаемое отдельно.

\textbf{Оценка $I_1$:}
$$ |I_1| = \left| \int_a^b (f(x) - g(x)) \sin(\omega x + \alpha) \, dx \right| \le \int_a^b |f(x) - g(x)| \cdot \underbrace{|\sin(\dots)|}_{\le 1} \, dx $$
$$ \le \int_a^b |f(x) - g(x)| \, dx = \| f - g \|_{L_1} < \frac{\varepsilon}{2} $$
(Это неравенство выполняется всегда, независимо от $\omega$, благодаря выбору функции $g$).

\textbf{Оценка $I_2$:}
Так как $g \in C^1([a, b])$, то мы можем применить доказанный пункт 1. Согласно пункту 1:
$$ \int_a^b g(x) \sin(\omega x + \alpha) \, dx \xrightarrow[\omega \to \infty]{} 0 $$
Это значит, что для выбранного $\varepsilon$ существует $N > 0$, такое что для всех $\omega > N$:
$$ |I_2| = \left| \int_a^b g(x) \sin(\omega x + \alpha) \, dx \right| < \frac{\varepsilon}{2} $$

\textbf{Итог для любого $\varepsilon > 0$:}
Для всех $\omega > N$ выполняется:
$$ \left| \int_a^b f(x) \sin(\omega x + \alpha) \, dx \right| \le |I_1| + |I_2| < \frac{\varepsilon}{2} + \frac{\varepsilon}{2} = \varepsilon $$

Так как $\varepsilon$ произвольно, это означает, что предел равен 0. $\blacksquare$

\vspace{1cm}

\section{Следствие}

Пусть $f \in L_1([-\pi, \pi])$. Тогда коэффициенты Фурье $a_n, b_n$ стремятся к нулю при $n \to \infty$.

$$ a_n, b_n \xrightarrow{n \to \infty} 0 $$

\textbf{Доказательство:}
Коэффициенты Фурье определяются формулами:
$$ a_n = \frac{1}{\pi} \int_{-\pi}^{\pi} f(x) \cos(nx) \, dx $$
$$ b_n = \frac{1}{\pi} \int_{-\pi}^{\pi} f(x) \sin(nx) \, dx $$

По доказанной Лемме Римана (полагая интервал $[a, b] = [-\pi, \pi]$, $\omega = n$, $\alpha = 0$ и $\alpha = \pi/2$ для косинуса):
$$ \int_{-\pi}^{\pi} f(x) \cos(nx) \, dx \xrightarrow{n \to \infty} 0 $$
$$ \int_{-\pi}^{\pi} f(x) \sin(nx) \, dx \xrightarrow{n \to \infty} 0 $$
Следовательно, $a_n \to 0$ и $b_n \to 0$.


\chapter{Билет 8 }
\section{Ядра Дирихле. Признак Дини и его обобщение на случай скачка функции}
\section{Ядро Дирихле}

\subsection{Определение и вывод формулы}
Пусть $f \in L_1([-\pi, \pi])$. Частичная сумма ряда Фурье $S_n(x)$ определяется как:
$$ S_n(x) = \sum_{k=-n}^{n} c_k e^{ikx}, $$
где коэффициенты Фурье $c_k$ вычисляются по формуле:
$$ c_k = \frac{1}{2\pi} \int_{-\pi}^{\pi} f(t) e^{-ikt} \, dt. $$

Подставим выражение для коэффициентов в сумму:
$$ S_n(x) = \sum_{k=-n}^{n} \left( \frac{1}{2\pi} \int_{-\pi}^{\pi} f(t) e^{-ikt} \, dt \right) e^{ikx} = \int_{-\pi}^{\pi} f(t) \left( \frac{1}{2\pi} \sum_{k=-n}^{n} e^{ik(x-t)} \right) \, dt. $$
Сделаем замену переменной $y = x - t$ (учитывая периодичность):
$$ S_n(x) = \int_{-\pi}^{\pi} f(x-t) D_n(t) \, dt. $$

Здесь $D_n(y)$ — это \textbf{Ядро Дирихле}:
$$ D_n(y) = \frac{1}{2\pi} \sum_{k=-n}^{n} e^{iky}. $$

Вычислим явный вид ядра Дирихле, используя формулу суммы геометрической прогрессии.
Сумма $\sum_{k=-n}^{n} e^{iky}$ представляет собой геометрическую прогрессию со знаменателем $q = e^{iy}$.
$$ \sum_{k=-n}^{n} e^{iky} = e^{-iny} + \dots + e^{iny} = e^{-iny} \frac{(e^{iy})^{2n+1} - 1}{e^{iy} - 1} = \frac{e^{i(n+1)y} - e^{-iny}}{e^{iy} - 1}. $$

Чтобы привести это к вещественному виду, вынесем половинные углы $e^{iy/2}$ в числителе и знаменателе:
$$ \frac{e^{i(n+1)y} - e^{-iny}}{e^{iy} - 1} = \frac{e^{iy/2} \left( e^{i(n+\frac{1}{2})y} - e^{-i(n+\frac{1}{2})y} \right)}{e^{iy/2} \left( e^{iy/2} - e^{-iy/2} \right)}. $$
Используя формулу Эйлера ($e^{i\phi} - e^{-i\phi} = 2i\sin\phi$), получаем:
$$ = \frac{2i \sin\left( \left(n + \frac{1}{2}\right)y \right)}{2i \sin\left(\frac{y}{2}\right)} = \frac{\sin\left( \left(n + \frac{1}{2}\right)y \right)}{\sin\left(\frac{y}{2}\right)}. $$

Таким образом, окончательная формула для ядра Дирихле:
$$ D_n(y) = \frac{1}{2\pi} \frac{\sin\left( \left(n + \frac{1}{2}\right)y \right)}{\sin\left(\frac{y}{2}\right)}. $$

\begin{center}
\begin{tikzpicture}
\begin{axis}[
    width=12cm, height=5cm,
    axis lines=middle,
    xmin=-7, xmax=7,
    ymin=-0.5, ymax=2,
    xlabel=$y$, ylabel=$D_n(y)$,
    title={График ядра Дирихле при $n=3$},
    xtick={-3.14, 0, 3.14},
    xticklabels={$-\pi$, $0$, $\pi$},
    samples=200, domain=-6.5:6.5
]
\addplot[blue, thick] {sin(deg((3+0.5)*x)) / (2*3.14159 * sin(deg(0.5*x)))};
\end{axis}
\end{tikzpicture}
\end{center}

\subsection{Свойства ядра Дирихле}
1.  **Нормировка:** Интеграл от ядра по периоду равен единице:
    $$ \int_{-\pi}^{\pi} D_n(x) \, dx = \frac{1}{2\pi} \sum_{k=-n}^{n} \int_{-\pi}^{\pi} e^{ikx} \, dx = 1, $$
    так как интеграл от $e^{ikx}$ равен 0 при $k \neq 0$ и $2\pi$ при $k=0$.
2.  **Неограниченность нормы:** Интеграл от модуля ядра (константы Лебега) расходится:
    $$ \int_{-\pi}^{\pi} |D_n(x)| \, dx \to \infty \quad \text{при} \quad n \to \infty. $$

\section{Признак Дини}

\textbf{Теорема (Признак Дини):}
Пусть $f \in L_1([-\pi, \pi])$.
Рассмотрим точку $x \in [-\pi, \pi]$, в которой функция определена.
Если существует $\delta > 0$ такое, что сходится интеграл:
$$ \int_{-\delta}^{\delta} \left| \frac{f(x+t) - f(x)}{t} \right| \, dt < \infty, $$
то частичные суммы ряда Фурье сходятся к значению функции в этой точке:
$$ S_n(x) \xrightarrow{n \to \infty} f(x). $$

\subsection{Доказательство}

Рассмотрим разность между частичной суммой и значением функции $f(x)$:
$$ S_n(x) - f(x) = \int_{-\pi}^{\pi} f(x-t) D_n(t) \, dt - f(x). $$
Используя свойство нормировки $\int_{-\pi}^{\pi} D_n(t) \, dt = 1$, представим $f(x)$ как:
$$ f(x) = f(x) \cdot 1 = \int_{-\pi}^{\pi} f(x) D_n(t) \, dt. $$
Тогда:
$$ S_n(x) - f(x) = \int_{-\pi}^{\pi} \big( f(x-t) - f(x) \big) D_n(t) \, dt. $$
Сделаем замену переменной $t \to -t$ в интеграле. Так как ядро Дирихле $D_n(t)$ — четная функция ($D_n(-t) = D_n(t)$), получаем:
$$ S_n(x) - f(x) = \int_{-\pi}^{\pi} \big( f(x+t) - f(x) \big) D_n(t) \, dt. $$

Подставим явное выражение для ядра Дирихле:
$$ S_n(x) - f(x) = \int_{-\pi}^{\pi} \frac{f(x+t) - f(x)}{2\pi \sin(t/2)} \sin\left(\left(n + \frac{1}{2}\right)t\right) \, dt. $$
Перепишем подынтегральное выражение в виде:
$$ = \int_{-\pi}^{\pi} \underbrace{\frac{f(x+t) - f(x)}{t} \cdot \frac{t}{2\pi \sin(t/2)}}_{g(t)} \cdot \sin\left(\left(n + \frac{1}{2}\right)t\right) \, dt. $$

Задача сводится к применению **Леммы Римана**. Если функция $g(t) \in L_1([-\pi, \pi])$ (абсолютно интегрируема), то по лемме Римана интеграл $\int g(t) \sin(\omega t) \, dt \to 0$ при $\omega \to \infty$.

\textbf{Проверка абсолютной интегрируемости $g(t)$:}
Нам нужно доказать, что $\int_{-\pi}^{\pi} |g(t)| \, dt < \infty$.
$$ |g(t)| = \left| \frac{f(x+t) - f(x)}{t} \right| \cdot \left| \frac{t}{2\pi \sin(t/2)} \right|. $$

Оценим множитель $\left| \frac{t}{2\pi \sin(t/2)} \right|$.
Известно неравенство Жордана: $\sin x \ge \frac{2}{\pi}x$ для $x \in [0, \pi/2]$.
Положим $x = |t|/2$. Так как $t \in [-\pi, \pi]$, то $|t|/2 \in [0, \pi/2]$.
$$ \sin\frac{|t|}{2} \ge \frac{2}{\pi} \cdot \frac{|t|}{2} = \frac{|t|}{\pi} \implies \frac{1}{\sin(|t|/2)} \le \frac{\pi}{|t|}. $$
Следовательно:
$$ \left| \frac{t}{2\pi \sin(t/2)} \right| = \frac{|t|}{2\pi} \cdot \frac{1}{|\sin(t/2)|} \le \frac{|t|}{2\pi} \cdot \frac{\pi}{|t|} = \frac{1}{2}. $$

Теперь оценим интеграл $I = \int_{-\pi}^{\pi} |g(t)| \, dt$. Разобьем его на три части: центральную окрестность $[-\delta, \delta]$ и два "хвоста" $[-\pi, -\delta]$ и $[\delta, \pi]$.
$$ I \le \frac{1}{2} \int_{-\pi}^{\pi} \left| \frac{f(x+t) - f(x)}{t} \right| \, dt = \frac{1}{2} \left( \int_{-\delta}^{\delta} + \int_{-\pi}^{-\delta} + \int_{\delta}^{\pi} \right) \left| \frac{f(x+t) - f(x)}{t} \right| \, dt. $$

1.  **Центральная часть** ($\int_{-\delta}^{\delta}$):
    Обозначим $I_1 = \int_{-\delta}^{\delta} \frac{|f(x+t) - f(x)|}{|t|} \, dt$. Этот интеграл сходится ($I_1 < \infty$) \textbf{по условию теоремы Дини}.

2.  **Хвосты** ($|t| \ge \delta$):
    На множестве $[-\pi, -\delta] \cup [\delta, \pi]$ выполняется неравенство $\frac{1}{|t|} \le \frac{1}{\delta}$.
    $$ \int_{\text{хвосты}} \frac{|f(x+t) - f(x)|}{|t|} \, dt \le \frac{1}{\delta} \int_{\text{хвосты}} (|f(x+t)| + |f(x)|) \, dt. $$
    Так как $f \in L_1$, то $\int |f(x+t)| \, dt \le \|f\|_{L_1}$. Величина $|f(x)|$ — константа.
    $$ \le \frac{1}{\delta} \left( \|f\|_{L_1} + 2\pi |f(x)| \right) < \infty. $$

\textbf{Вывод:} Функция $g(t)$ абсолютно интегрируема ($g \in L_1$). Следовательно, применима лемма Римана:
$$ \int_{-\pi}^{\pi} g(t) \sin\left(\left(n + \frac{1}{2}\right)t\right) \, dt \xrightarrow{n \to \infty} 0. $$
То есть $S_n(x) \to f(x)$. $\blacksquare$

\begin{center}
\begin{tikzpicture}
\draw[->] (-4,0) -- (4,0) node[right] {$t$};
\draw[thick] (-3.5,0) -- (-3.5, 0.2); \node[below] at (-3.5,0) {$-\pi$};
\draw[thick] (3.5,0) -- (3.5, 0.2); \node[below] at (3.5,0) {$\pi$};
\draw[thick] (-1,0) -- (-1, 0.2); \node[below] at (-1,0) {$-\delta$};
\draw[thick] (1,0) -- (1, 0.2); \node[below] at (1,0) {$\delta$};
\draw[thick] (0,0) -- (0, 0.2); \node[below] at (0,0) {$0$};

\draw[<->, blue] (-1, 0.5) -- (1, 0.5) node[midway, above] {Условие Дини};
\draw[<->, red] (1, 0.5) -- (3.5, 0.5) node[midway, above] {$f \in L_1$};
\draw[<->, red] (-3.5, 0.5) -- (-1, 0.5) node[midway, above] {$f \in L_1$};
\end{tikzpicture}
\end{center}

\subsection{Замечание: достаточные условия}
Условие Дини выполняется, если:
\begin{enumerate}
    \item Существует производная $f'(x)$. Тогда $\lim_{t \to 0} \frac{f(x+t)-f(x)}{t} = f'(x)$, функция ограничена, интеграл сходится.
    \item Существуют односторонние производные $f'_{\pm}(x)$ (функция дифференцируема слева и справа).
    \item Выполняется \textbf{условие Гёльдера}: $|f(x+t) - f(x)| \le C |t|^\alpha$ при $|t| < \delta$ и $\alpha \in (0, 1]$.
    Действительно, $\int \frac{|t|^\alpha}{|t|} dt = \int |t|^{\alpha-1} dt$ сходится, так как $\alpha - 1 > -1$.
\end{enumerate}

\section{Обобщение на случай скачка функции}

\textbf{Теорема:}
Пусть $f \in L_1([-\pi, \pi])$ и имеет в точке $x$ разрыв первого рода (существуют пределы слева $f(x-0)$ и справа $f(x+0)$).
Рассмотрим пределы:
$$ \lim_{t \to 0-} f(x+t) = f(x-0), \quad \lim_{t \to 0+} f(x+t) = f(x+0). $$
Если существует $\delta > 0$, такое что сходятся интегралы (обобщенное условие Дини):
$$ \int_{-\delta}^{0} \left| \frac{f(x+t) - f(x-0)}{t} \right| \, dt < \infty \quad \text{и} \quad \int_{0}^{\delta} \left| \frac{f(x+t) - f(x+0)}{t} \right| \, dt < \infty, $$
то частичные суммы ряда Фурье сходятся к полусумме пределов:
$$ S_n(x) \xrightarrow{n \to \infty} \frac{f(x-0) + f(x+0)}{2}. $$

\subsection{Доказательство}

Используем свойство ядра Дирихле. Так как $D_n(t)$ четная и $\int_{-\pi}^{\pi} D_n(t) dt = 1$, то:
$$ \int_{-\pi}^{0} D_n(t) \, dt = \frac{1}{2}, \quad \int_{0}^{\pi} D_n(t) \, dt = \frac{1}{2}. $$

Запишем выражение для предельных значений:
$$ \frac{f(x-0)}{2} = f(x-0) \int_{-\pi}^{0} D_n(t) \, dt = \int_{-\pi}^{0} f(x-0) D_n(t) \, dt, $$
$$ \frac{f(x+0)}{2} = f(x+0) \int_{0}^{\pi} D_n(t) \, dt = \int_{0}^{\pi} f(x+0) D_n(t) \, dt. $$

Рассмотрим разность между $S_n(x)$ и целевым пределом. Разобьем интеграл $S_n(x)$ на два промежутка: $[-\pi, 0]$ и $[0, \pi]$.
\begin{align*}
S_n(x) - \frac{f(x-0) + f(x+0)}{2} &= \left( \int_{-\pi}^{0} f(x+t)D_n(t)dt + \int_{0}^{\pi} f(x+t)D_n(t)dt \right) \\
&- \left( \int_{-\pi}^{0} f(x-0)D_n(t)dt + \int_{0}^{\pi} f(x+0)D_n(t)dt \right).
\end{align*}

Сгруппируем слагаемые:
$$ = \underbrace{\int_{-\pi}^{0} \big( f(x+t) - f(x-0) \big) D_n(t) \, dt}_{I_1} + \underbrace{\int_{0}^{\pi} \big( f(x+t) - f(x+0) \big) D_n(t) \, dt}_{I_2}. $$

Рассмотрим, например, второй интеграл $I_2$ (для $t > 0$):
$$ I_2 = \int_{0}^{\pi} \frac{f(x+t) - f(x+0)}{2\pi \sin(t/2)} \sin\left(\left(n + \frac{1}{2}\right)t\right) \, dt. $$
При малых $t$ имеем $\sin(t/2) \sim t/2$.
Сходимость $I_2$ к нулю при $n \to \infty$ зависит от абсолютной интегрируемости функции:
$$ g(t) = \frac{f(x+t) - f(x+0)}{t}. $$
Это и есть условие теоремы (для правого интервала). Аналогично для $I_1$ и левого интервала.

Так как оба интеграла стремятся к нулю (по лемме Римана), получаем:
$$ S_n(x) \to \frac{f(x-0) + f(x+0)}{2}. \quad \blacksquare $$

\chapter{Билет 9}
\section{Билет: Принцип локализации Римана. Поведение коэффициентов Фурье при n стремящейся к бесконечности.}

\subsection{Принцип локализации Римана}

\textbf{Утверждение:}
Пусть $f \in L_1([-\pi, \pi])$.
На сходимость ряда Фурье $S_n(x) \to f(x)$ в данной точке $x$ \textbf{не влияют} значения функции вне любой сколь угодно малой окрестности $(x-\delta, x+\delta)$ при $\forall \delta > 0$.

То есть поведение ряда Фурье в точке $x$ зависит только от поведения функции в бесконечно малой окрестности этой точки.

\subsubsection*{Доказательство}

1. Запишем частичную сумму ряда Фурье $S_n(x)$ в интегральной форме с использованием ядра Дирихле:
$$ S_n(x) = \int_{-\pi}^{\pi} f(x+t) D_n(t) \, dt, $$
где $D_n(t)$ — ядро Дирихле:
$$ D_n(t) = \frac{\sin\left((n+\frac{1}{2})t\right)}{2\pi \sin(t/2)}. $$

2. Разобьем интеграл по отрезку $[-\pi, \pi]$ на три части: интеграл по окрестности нуля $[-\delta, \delta]$ и два интеграла по «хвостам» (где $|t| \ge \delta$):
$$ S_n(x) = \left( \int_{-\pi}^{-\delta} + \int_{\delta}^{\pi} \right) f(x+t) D_n(t) \, dt + \int_{-\delta}^{\delta} f(x+t) D_n(t) \, dt. $$

3. Рассмотрим интегралы по внешним областям (хвостам), где $|t| \ge \delta$.
Заметим, что если $|t| \in [\delta, \pi]$, то знаменатель ядра Дирихле отделен от нуля:
$$ \sin(t/2) \neq 0. $$

4. Введем вспомогательную функцию $g(t)$:
$$ g(t) = \frac{f(x+t)}{2\pi \sin(t/2)}. $$
Так как $f \in L_1([-\pi, \pi])$ и функция $\frac{1}{\sin(t/2)}$ ограничена на множестве $[-\pi, \pi] \setminus [-\delta, \delta]$, то функция $g(t)$ является абсолютно интегрируемой на этом множестве:
$$ g(t) \in L_1\left([-\pi, \pi] \setminus [-\delta, \delta]\right). $$

5. Тогда сумму интегралов по хвостам можно переписать как:
$$ \int_{\text{хвосты}} g(t) \sin\left(\left(n+\frac{1}{2}\right)t\right) \, dt. $$
По \textbf{лемме Римана}, этот интеграл стремится к нулю при $n \to \infty$:
$$ \int_{\text{хвосты}} g(t) \sin\left(\left(n+\frac{1}{2}\right)t\right) \, dt \xrightarrow{n \to \infty} 0. $$

6. Следовательно:
$$ S_n(x) = o(1) + \int_{-\delta}^{\delta} f(x+t) D_n(t) \, dt. $$
Предел $S_n(x)$ полностью определяется вторым слагаемым — интегралом по малой окрестности $(-\delta, \delta)$. $\blacksquare$

\vspace{0.5cm}
\begin{center}
\begin{tikzpicture}
    % Ось
    \draw[->] (-5,0) -- (5,0) node[right] {$t$};
    
    % Границы
    \draw[thick] (-4,0.2) -- (-4,-0.2) node[below] {$-\pi$};
    \draw[thick] (4,0.2) -- (4,-0.2) node[below] {$\pi$};
    \draw[thick] (-1.5,0.2) -- (-1.5,-0.2) node[below] {$-\delta$};
    \draw[thick] (1.5,0.2) -- (1.5,-0.2) node[below] {$\delta$};
    \draw (0,0.1) -- (0,-0.1) node[below] {$0$};
    
    % Области
    \draw[<->, red, thick] (-4, 0.5) -- (-1.5, 0.5) node[midway, above] {Вклад $\to 0$};
    \draw[<->, blue, thick] (-1.5, 0.5) -- (1.5, 0.5) node[midway, above] {Определяет сходимость};
    \draw[<->, red, thick] (1.5, 0.5) -- (4, 0.5) node[midway, above] {Вклад $\to 0$};
\end{tikzpicture}
\end{center}

\newpage

\subsection{Поведение коэффициентов Фурье гладкой функции}

\textbf{Теорема (о порядке убывания коэффициентов):}

Пусть $f \in C^k([-\pi, \pi])$ (функция имеет $k$ непрерывных производных) и выполнены условия периодичности для всех производных до $(k-1)$-го порядка включительно:
$$ f^{(j)}(-\pi) = f^{(j)}(\pi) \quad \forall j = 0, 1, \dots, k-1. $$

Тогда коэффициенты Фурье $a_n, b_n$ убывают быстрее, чем $\frac{1}{n^k}$ при $n \to \infty$:
$$ a_n, b_n = o\left(\frac{1}{n^k}\right). $$

\subsubsection*{Доказательство}

Докажем утверждение для коэффициентов $b_n$ (для $a_n$ доказательство аналогично).

1. Рассмотрим формулу для коэффициента $b_n$:
$$ b_n = \frac{1}{\pi} \int_{-\pi}^{\pi} f(x) \sin(nx) \, dx. $$

2. Применим метод \textbf{интегрирования по частям}:
\begin{align*}
u &= f(x) \quad \Rightarrow \quad du = f'(x) \, dx \\
dv &= \sin(nx) \, dx \quad \Rightarrow \quad v = -\frac{1}{n} \cos(nx)
\end{align*}

Подставим в формулу интегрирования по частям:
$$ b_n = \frac{1}{\pi} \left( \left[ -\frac{f(x)}{n} \cos(nx) \right]_{-\pi}^{\pi} + \frac{1}{n} \int_{-\pi}^{\pi} f'(x) \cos(nx) \, dx \right). $$

3. Рассмотрим граничное слагаемое (внеинтегральный член):
$$ -\frac{1}{\pi n} \left( f(\pi)\cos(n\pi) - f(-\pi)\cos(-n\pi) \right). $$
Так как $\cos(n\pi) = (-1)^n$ и $\cos(-n\pi) = (-1)^n$, а по условию теоремы $f(\pi) = f(-\pi)$, то:
$$ f(\pi)(-1)^n - f(-\pi)(-1)^n = 0. $$
Слагаемое зануляется.

4. Получаем выражение после первого шага:
$$ b_n = \frac{1}{n} \cdot \underbrace{\frac{1}{\pi} \int_{-\pi}^{\pi} f'(x) \cos(nx) \, dx}_{\text{связано с коэф. Фурье для } f'}. $$

5. \textbf{Повторим интегрирование по частям $k$ раз.}
Каждый раз будет выноситься множитель вида $\frac{1}{n}$ или $-\frac{1}{n}$.
Граничные члены будут зануляться из-за условий периодичности производных ($f^{(j)}(\pi) = f^{(j)}(-\pi)$).

В итоге получим:
$$ b_n = \pm \frac{1}{n^k} \cdot \frac{1}{\pi} \int_{-\pi}^{\pi} f^{(k)}(x) \cdot \left\{ \begin{matrix} \sin(nx) \\ \cos(nx) \end{matrix} \right\} \, dx. $$

6. Так как по условию $f^{(k)}$ непрерывна (а значит, интегрируема), то к интегралу
$$ \int_{-\pi}^{\pi} f^{(k)}(x) \cdot \dots \, dx $$
применимо \textbf{Следствие из леммы Римана}: коэффициенты Фурье интегрируемой функции стремятся к нулю.
Обозначим этот интеграл через $\varepsilon_n$, где $\varepsilon_n \xrightarrow{n \to \infty} 0$.

7. Тогда:
$$ b_n = \frac{1}{n^k} \cdot \varepsilon_n = o\left(\frac{1}{n^k}\right). \quad \blacksquare $$

\subsection*{Пример ($k=2$)}
Пусть $f \in C^2([-\pi, \pi])$ и выполнены условия периодичности:
$$ f(-\pi) = f(\pi), \quad f'(-\pi) = f'(\pi). $$
Тогда коэффициенты Фурье убывают как "o-малое" от обратного квадрата:

$$ a_n, b_n = o\left(\frac{1}{n^2}\right). $$

\chapter{Билет 10}
\section{Теорема о $\delta$-образной последовательности}

Пусть $a < 0$, $b > 0$. Рассмотрим интервал $[a, b]$.

\subsection{Определение}
Последовательность функций $\{g_n\}_{n=1}^{\infty}$ из пространства $L_1([a, b])$ называется \textbf{дельтаобразной} (или $\delta$-образной), если выполнены три условия:

\begin{enumerate}
    \item \textbf{Интеграл нормирован:}
    \[
    \int\limits_{a}^{b} g_n(t)\,dt \xrightarrow{n \to \infty} 1
    \]
    
    \item \textbf{Равномерная ограниченность норм:}
    Существует константа $M > 0$, такая, что для всех $n \in \mathbb{N}$:
    \[
    \int\limits_{a}^{b} |g_n(t)|\,dt < M
    \]
    
    \item \textbf{Стягивание к нулю (концентрация массы):}
    Для любого $\delta > 0$ интеграл вне $\delta$-окрестности нуля стремится к нулю:
    \[
    \left( \int\limits_{a}^{-\delta} + \int\limits_{\delta}^{b} \right) |g_n(t)|\,dt \xrightarrow{n \to \infty} 0
    \]
\end{enumerate}

\begin{center}
\begin{tikzpicture}
    % Axes
    \draw[->] (-4,0) -- (4,0) node[right] {$t$};
    \draw[->] (0,-0.5) -- (0,4) node[above] {$g_n(t)$};
    
    % Labels
    \node at (0,0) [below left] {$0$};
    \node at (-3.5,0) [below] {$a$};
    \node at (3.5,0) [below] {$b$};
    \node at (-1,0) [below] {$-\delta$};
    \node at (1,0) [below] {$\delta$};
    
    % Function n=1 (Wide)
    \draw[blue, thick] plot [smooth] coordinates {(-3.5,0.1) (-2,0.2) (0,1) (2,0.2) (3.5,0.1)};
    \node[blue] at (2,0.8) {$n=1$};
    
    % Function n=large (Narrow)
    \draw[red, thick] plot [smooth] coordinates {(-3.5,0.02) (-1,0.05) (-0.5, 0.5) (0,3.5) (0.5, 0.5) (1,0.05) (3.5,0.02)};
    \node[red] at (0.8,3) {$n \to \infty$};
    
    % Highlight tails
    \draw[dashed] (-1,0) -- (-1, 1);
    \draw[dashed] (1,0) -- (1, 1);
    
    \node at (2.5, -0.7) {\small Основная масса стягивается к 0};
\end{tikzpicture}
\end{center}

\subsection{Теорема}
Если $\{g_n\}_{n=1}^{\infty}$ — $\delta$-образная последовательность на $[a, b]$, то для любой непрерывной функции $f \in C([a, b])$ выполняется предельное соотношение:
\[
\int\limits_{a}^{b} g_n(t) f(t)\,dt \xrightarrow{n \to \infty} f(0)
\]

\subsection{Доказательство}

\textbf{Шаг 1. Предварительные преобразования.}
Достаточно доказать, что
\[
\int\limits_{a}^{b} g_n(t) (f(t) - f(0))\,dt \xrightarrow{n \to \infty} 0
\]
Действительно, мы можем записать исходный интеграл следующим образом:
\[
\int\limits_{a}^{b} g_n f\,dt = \int\limits_{a}^{b} g_n(f - f(0))\,dt + f(0)\int\limits_{a}^{b} g_n(t)\,dt
\]
Второе слагаемое $f(0)\int_{a}^{b} g_n(t)\,dt \to f(0) \cdot 1 = f(0)$ в силу условия 1 (нормировка интеграла). Следовательно, задача сводится к доказательству того, что первое слагаемое стремится к нулю.

\textbf{Шаг 2. Выбор окрестности (использование непрерывности).}
Так как функция $f$ непрерывна в точке 0, то для любого $\varepsilon > 0$ существует $\delta < \min\{|a|, b\}$, такое что для всех $t \in (-\delta, \delta)$:
\[
|f(t) - f(0)| < \frac{\varepsilon}{2M},
\]
где $M$ — константа из условия 2.

\textbf{Шаг 3. Оценка функции на всем отрезке.}
Так как функция $f$ непрерывна на компакте $[a, b]$, то она ограничена. Пусть норма в пространстве непрерывных функций определена как $\|f\|_{C[a,b]} = \max_{t \in [a,b]} |f(t)|$.
Тогда для любого $t \in [a, b]$ справедливо неравенство треугольника:
\[
|f(t) - f(0)| \le |f(t)| + |f(0)| \le 2\|f\|_{C[a,b]}.
\]

По условию 3 (стягивание к нулю), для выбранного $\delta$ и любого $\varepsilon$ существует номер $N$, такой что для всех $n > N$:
\[
\left( \int\limits_{a}^{-\delta} + \int\limits_{\delta}^{b} \right) |g_n(t)|\,dt < \frac{\varepsilon}{4\|f\|_{C[a,b]}}.
\]

\textbf{Шаг 4. Разбиение интеграла и итоговая оценка.}
Рассмотрим модуль интеграла разности для всех $n > N$. Разобьем область интегрирования на две части: внешнюю часть («хвосты») и окрестность нуля.
\[
\left| \int\limits_{a}^{b} g_n(t)(f(t) - f(0))\,dt \right| \le 
\underbrace{\left( \int\limits_{a}^{-\delta} + \int\limits_{\delta}^{b} \right) |g_n(t)| \cdot |f(t) - f(0)|\,dt}_{\text{Часть 1}} + 
\underbrace{\int\limits_{-\delta}^{\delta} |g_n(t)| \cdot |f(t) - f(0)|\,dt}_{\text{Часть 2}}
\]

\textbf{Оценим Часть 1 (интеграл по «хвостам»):}
Используем глобальную оценку $|f(t) - f(0)| \le 2\|f\|_C$:
\[
\left( \int\limits_{a}^{-\delta} + \int\limits_{\delta}^{b} \right) |g_n(t)| \cdot |f(t) - f(0)|\,dt \le 2\|f\|_C \cdot \left( \int\limits_{a}^{-\delta} + \int\limits_{\delta}^{b} \right) |g_n(t)|\,dt
\]
Используя следствие из условия 3 для $n > N$:
\[
< 2\|f\|_C \cdot \frac{\varepsilon}{4\|f\|_C} = \frac{\varepsilon}{2}.
\]

\textbf{Оценим Часть 2 (интеграл в окрестности нуля):}
Внутри интервала $(-\delta, \delta)$ известно, что $|f(t) - f(0)| < \frac{\varepsilon}{2M}$.
\[
\int\limits_{-\delta}^{\delta} |g_n(t)| \cdot |f(t) - f(0)|\,dt \le \frac{\varepsilon}{2M} \int\limits_{-\delta}^{\delta} |g_n(t)|\,dt
\]
Расширим интервал интегрирования до всего отрезка $[a, b]$ (так как $|g_n| \ge 0$, интеграл может только увеличиться) и воспользуемся условием 2:
\[
\le \frac{\varepsilon}{2M} \int\limits_{a}^{b} |g_n(t)|\,dt < \frac{\varepsilon}{2M} \cdot M = \frac{\varepsilon}{2}.
\]

\textbf{Заключение:}
Складывая оценки двух частей, получаем:
\[
\left| \int\limits_{a}^{b} g_n(t)(f(t) - f(0))\,dt \right| < \frac{\varepsilon}{2} + \frac{\varepsilon}{2} = \varepsilon.
\]
Что и требовалось доказать. $\blacksquare$

\subsection{Замечания}

\textbf{Замечание 1:}
Условие 2 выполняется \textbf{автоматически}, если $g_n(t) \ge 0$ для всех $n$.
Так как если $f$ неотрицательна, то $\int |g_n| = \int g_n$. Из условия 1 следует, что $\int g_n \to 1$, следовательно, последовательность интегралов ограничена (сходящаяся числовая последовательность всегда ограничена).

\textbf{Замечание 2 (Связь с функциональным анализом):}
Данная теорема означает сходимость в пространстве $C^*([a, b])$.
$C^*([a, b])$ — это пространство комплекснозначных мер на $[a, b]$, сопряженное к пространству непрерывных функций $C([a, b])$. Предел является $\delta$-функцией Дирака.

Функционал $\delta$ действует на $f \in C([a, b])$ следующим образом: $(\delta, f) = f(0)$.
Это соответствует мере Стилтьеса, порожденной функцией Хевисайда $\theta(x)$:
\[
\theta(x) = \begin{cases} 1, & x \ge 0 \\ 0, & x < 0 \end{cases} \quad \text{(или для множеств } E \in \mathcal{B}: \theta(E) = \begin{cases} 1, & 0 \in E \\ 0, & 0 \notin E \end{cases} \text{)}
\]
Действие функционала $\mu$ на функцию $f$ записывается как:
\[
(\mu, f) = \int\limits_{a}^{b} f\,d\mu, \quad |(\mu, f)| \le \|f\|_{C[a,b]} \cdot |\mu| \quad \text{($|\mu|$ --- полная вариация меры)}.
\]
Такая сходимость называется \textbf{*-слабой} (звездочка-слабой) сходимостью:
\[
g_n(x)\,dx \xrightarrow[n \to \infty]{*-w} \delta(x) \quad \text{в } C^*([a, b])
\]
Это означает, что для любой $f \in C([a, b])$:
\[
\int\limits_{a}^{b} f(x) g_n(x)\,dx \xrightarrow{n \to \infty} \int\limits_{a}^{b} f(x)\,d\mu_{\delta}(x) = (\delta, f) = f(0).
\]

\chapter{Билет 11}

\section{ Ядра Фейера. Суммирование рядов по Чезаро. Теорема Фейера}

\subsection{1. Постановка задачи и определение средних Чезаро}

Пусть функция $f \in L_1([-\pi, \pi])$. Ряд Фурье функции $f$ имеет вид:
$$ \frac{a_0}{2} + \sum_{k=1}^{\infty} (a_k \cos(kx) + b_k \sin(kx)) $$

Обозначим через $S_n(x)$ частичную сумму ряда Фурье:
$$ S_n(x) = \frac{a_0}{2} + \sum_{k=1}^{n} (a_k \cos(kx) + b_k \sin(kx)) $$

Известно, что последовательность частичных сумм $S_n(x)$ может расходиться даже для непрерывных функций. Поэтому вместо сходимости сумм рассматривают сходимость их средних арифметических (метод суммирования Чезаро).

\textbf{Определение (Суммы Фейера/Средние Чезаро).} \\
Суммами Фейера $\sigma_n(x)$ называется среднее арифметическое первых $n$ частичных сумм ряда Фурье:
\begin{equation}
    \sigma_n(x) = \frac{S_0(x) + S_1(x) + \dots + S_{n-1}(x)}{n}
\end{equation}

Так как частичная сумма представима через интеграл с ядром Дирихле $S_k(x) = \int_{-\pi}^{\pi} f(x-t)D_k(t)\,dt$, то, подставляя это в определение $\sigma_n(x)$ и пользуясь линейностью интеграла, получаем:
$$ \sigma_n(x) = \frac{1}{n} \sum_{k=0}^{n-1} \int_{-\pi}^{\pi} f(x-t) D_k(t)\,dt = \int_{-\pi}^{\pi} f(x-t) \left[ \frac{1}{n} \sum_{k=0}^{n-1} D_k(t) \right] dt $$

\textbf{Определение (Ядро Фейера).} \\
Функция $\Phi_n(t)$, равная среднему арифметическому ядер Дирихле, называется ядром Фейера:
\begin{equation}
    \Phi_n(t) = \frac{1}{n} \sum_{k=0}^{n-1} D_k(t)
\end{equation}

---

\subsection{2. Явный вид ядра Фейера}

\textbf{Лемма.} Справедлива следующая формула для ядра Фейера на $[-\pi, \pi]$:
\begin{equation}
    \Phi_n(x) = \frac{1}{2\pi n} \frac{\sin^2\left(\frac{nx}{2}\right)}{\sin^2\left(\frac{x}{2}\right)}
\end{equation}

\textbf{Доказательство:} \\
Напомним формулу для ядра Дирихле:
$$ D_k(x) = \frac{1}{2\pi} \frac{\sin\left((k+\frac{1}{2})x\right)}{\sin\left(\frac{x}{2}\right)} $$
Нам необходимо найти сумму числителей. Рассмотрим эту сумму как мнимую часть суммы геометрической прогрессии.
$$ \sum_{k=0}^{n-1} \sin\left((k+\frac{1}{2})x\right) = \text{Im} \left( \sum_{k=0}^{n-1} e^{i(k+\frac{1}{2})x} \right) = \text{Im} \left( e^{i\frac{x}{2}} \sum_{k=0}^{n-1} (e^{ix})^k \right) $$
Сумма геометрической прогрессии равна $\frac{q^n - 1}{q - 1}$, где $q = e^{ix}$:
$$ = \text{Im} \left( e^{i\frac{x}{2}} \frac{e^{inx} - 1}{e^{ix} - 1} \right) $$
Преобразуем знаменатель, используя формулу Эйлера: $e^{ix} - 1 = e^{i\frac{x}{2}}(e^{i\frac{x}{2}} - e^{-i\frac{x}{2}}) = e^{i\frac{x}{2}} \cdot 2i \sin\frac{x}{2}$.
Подставляем:
$$ = \text{Im} \left( e^{i\frac{x}{2}} \frac{e^{inx} - 1}{e^{i\frac{x}{2}} \cdot 2i \sin\frac{x}{2}} \right) = \text{Im} \left( \frac{e^{inx} - 1}{2i \sin\frac{x}{2}} \right) $$
Выносим вещественный множитель $\frac{1}{2\sin(x/2)}$ за знак мнимой части. Учитывая, что $\frac{1}{i} = -i$:
$$ = \frac{1}{2\sin\frac{x}{2}} \text{Im} \left( -i (\cos nx + i\sin nx - 1) \right) $$
$$ = \frac{1}{2\sin\frac{x}{2}} \text{Im} \left( -i \cos nx + \sin nx + i \right) $$
Мнимая часть выражения равна $(1 - \cos nx)$. Используем формулу $1 - \cos nx = 2\sin^2\left(\frac{nx}{2}\right)$.
Итого сумма числителей равна $\frac{\sin^2(nx/2)}{\sin(x/2)}$.
Возвращаясь к формуле для $\Phi_n(x)$, добавляем множитель $\frac{1}{n}$ и знаменатель от $D_k$:
$$ \Phi_n(x) = \frac{1}{n} \cdot \frac{1}{2\pi \sin\frac{x}{2}} \cdot \frac{\sin^2\left(\frac{nx}{2}\right)}{\sin\frac{x}{2}} = \frac{1}{2\pi n} \frac{\sin^2\left(\frac{nx}{2}\right)}{\sin^2\left(\frac{x}{2}\right)} \quad \blacksquare $$

\begin{center}
\begin{tikzpicture}
\begin{axis}[
    width=12cm, height=7cm,
    domain=-3.14:3.14, samples=300,
    axis lines=middle,
    xlabel=$x$, ylabel=$\Phi_n(x)$,
    ymin=0, ymax=2.5,
    xtick={-3.14, 0, 3.14},
    xticklabels={$-\pi$, $0$, $\pi$},
    title={График ядра Фейера $\Phi_n(x)$ для $n=5$}
]
% График, избегая деления на ноль в точке 0
\addplot [blue, thick, domain=-3.14:-0.05] { (sin(deg(5*x/2))^2) / (5 * 3.14 * sin(deg(x/2))^2) };
\addplot [blue, thick, domain=0.05:3.14] { (sin(deg(5*x/2))^2) / (5 * 3.14 * sin(deg(x/2))^2) };
% Пик в нуле вручную, так как формула дает 0/0
\draw[blue, thick] (0, 2.5) -- (0, 1.6); % Схематично соединяем
\node at (axis cs: 1.5, 1.5) [anchor=west] {Основной пик растет $\sim n$};
\node at (axis cs: 1.5, 1.2) [anchor=west] {Хвосты убывают $\to 0$};
\end{axis}
\end{tikzpicture}
\end{center}

---

\subsection{3. Свойства ядра Фейера}

\textbf{Лемма.} Ядра Фейера образуют $\delta$-образную последовательность на $[-\pi, \pi]$.

Необходимо проверить три свойства:

\textbf{1. Нормировка.}
$$ \int_{-\pi}^{\pi} \Phi_n(x)\,dx = \frac{1}{n} \sum_{k=0}^{n-1} \int_{-\pi}^{\pi} D_k(x)\,dx = \frac{1}{n} \sum_{k=0}^{n-1} 1 = 1 $$

\textbf{2. Неотрицательность.}
Из формулы с квадратами синусов очевидно, что $\Phi_n(x) \ge 0$.
Отсюда автоматически следует равномерная ограниченность норм в $L_1$: $\int |\Phi_n| = \int \Phi_n = 1$.

\textbf{3. Стягивание к нулю (концентрация массы).}
Для любого фиксированного $\delta > 0$ рассмотрим интеграл по области вне окрестности нуля $[\delta, \pi]$.
На интервале $[\delta, \pi]$ знаменатель ограничен снизу: $\sin^2(x/2) \ge \sin^2(\delta/2)$.
Числитель всегда $\le 1$.
$$ \Phi_n(x) \le \frac{1}{2\pi n \sin^2(\delta/2)} $$
Следовательно:
$$ \int_{\delta \le |x| \le \pi} \Phi_n(x)\,dx \le \int_{\delta \le |x| \le \pi} \frac{C}{n}\,dx = O\left(\frac{1}{n}\right) \xrightarrow{n \to \infty} 0 $$

---

\subsection{4. Теорема Фейера}

\textbf{Теорема.} Пусть $f \in C([-\pi, \pi])$ и выполняется условие периодичности $f(-\pi) = f(\pi)$. Тогда средние Чезаро $\sigma_n(x)$ сходятся к функции $f(x)$ \textbf{равномерно} на отрезке $[-\pi, \pi]$:
$$ \sigma_n(x) \rightrightarrows f(x) \quad \text{при } n \to \infty $$

\textbf{Доказательство:} \\
1. Запишем разность, используя интегральное представление. Продолжим $f$ $2\pi$-периодически на $\mathbb{R}$.
$$ \sigma_n(x) - f(x) = \int_{-\pi}^{\pi} f(x-t)\Phi_n(t)\,dt - f(x) \cdot 1 $$
Так как $\int \Phi_n(t)dt = 1$, внесем $f(x)$ под интеграл:
$$ \sigma_n(x) - f(x) = \int_{-\pi}^{\pi} [f(x-t) - f(x)] \Phi_n(t)\,dt $$

2. Оценим модуль разности:
$$ |\sigma_n(x) - f(x)| \le \int_{-\pi}^{\pi} |f(x-t) - f(x)| \Phi_n(t)\,dt $$

3. Используем равномерную непрерывность.
Так как $f$ непрерывна на компакте и периодична, она равномерно непрерывна на всей оси.
$$ \forall \varepsilon > 0 \quad \exists \delta > 0: \quad \forall x, t \text{ таких, что } |t| < \delta \implies |f(x-t) - f(x)| < \frac{\varepsilon}{2} $$
\textit{Важно:} $\delta$ зависит только от $\varepsilon$, но не от $x$.

4. Разобьем интеграл на две части: $I_1$ (окрестность нуля) и $I_2$ (хвосты).
$$ I = \underbrace{\int_{-\delta}^{\delta} |f(x-t) - f(x)| \Phi_n(t)\,dt}_{I_1} + \underbrace{\left( \int_{-\pi}^{-\delta} + \int_{\delta}^{\pi} \right) |f(x-t) - f(x)| \Phi_n(t)\,dt}_{I_2} $$

\textbf{Оценка $I_1$:}
$$ I_1 < \int_{-\delta}^{\delta} \frac{\varepsilon}{2} \Phi_n(t)\,dt = \frac{\varepsilon}{2} \int_{-\delta}^{\delta} \Phi_n(t)\,dt \le \frac{\varepsilon}{2} \int_{-\pi}^{\pi} \Phi_n(t)\,dt = \frac{\varepsilon}{2} $$

\textbf{Оценка $I_2$:}
Функция $f$ ограничена, пусть $M = \sup |f(x)|$. Тогда $|f(x-t) - f(x)| \le 2M$.
$$ I_2 \le 2M \int_{\delta \le |t| \le \pi} \Phi_n(t)\,dt $$
По свойству ядра Фейера, интеграл по хвостам стремится к 0. Значит, $\exists N$, такое что для всех $n > N$:
$$ \int_{\delta \le |t| \le \pi} \Phi_n(t)\,dt < \frac{\varepsilon}{4M} \implies I_2 < \frac{\varepsilon}{2} $$

\textbf{Итог:}
Для любого $\varepsilon > 0$ при $n > N$ для всех $x \in [-\pi, \pi]$:
$$ |\sigma_n(x) - f(x)| \le I_1 + I_2 < \varepsilon $$
Что означает равномерную сходимость. $\blacksquare$
\chapter{Билет 12}
\section{ Интеграл Фурье в тригонометрической и комплексной формах (нестрогий вывод)}

\subsection{ Постановка задачи и условия}
Пусть функция $f$ удовлетворяет следующим условиям:
\begin{itemize}
    \item $f \in C(\mathbb{R}) \cap L_1(\mathbb{R})$ (функция непрерывна и абсолютно интегрируема на всей оси $\mathbb{R}$);
    \item В каждой точке $x \in \mathbb{R}$ выполнено условие Дини (или иное условие локальной сходимости ряда Фурье).
\end{itemize}

Рассмотрим разложение функции в ряд Фурье на симметричном отрезке $[-l, l]$ при $l \to \infty$.

\subsection{ Тригонометрическая форма интеграла Фурье}
Запишем ряд Фурье для функции $f(x)$ на промежутке $[-l, l]$:
\begin{equation}
    f(x) = \frac{a_0}{2} + \sum_{k=1}^{\infty} \left( a_k \cos\left(\frac{k\pi x}{l}\right) + b_k \sin\left(\frac{k\pi x}{l}\right) \right),
\end{equation}
где коэффициенты Фурье определяются формулами:
$$
a_k = \frac{1}{l} \int_{-l}^{l} f(t) \cos\left(\frac{k\pi t}{l}\right) dt, \quad k \ge 0,
$$
$$
b_k = \frac{1}{l} \int_{-l}^{l} f(t) \sin\left(\frac{k\pi t}{l}\right) dt, \quad k \ge 1.
$$

Подставим выражения для коэффициентов $a_k$ и $b_k$ обратно в ряд Фурье:
$$
f(x) = \frac{1}{2l} \int_{-l}^{l} f(t) dt + \sum_{k=1}^{\infty} \frac{1}{l} \int_{-l}^{l} f(t) \left[ \cos\left(\frac{k\pi t}{l}\right) \cos\left(\frac{k\pi x}{l}\right) + \sin\left(\frac{k\pi t}{l}\right) \sin\left(\frac{k\pi x}{l}\right) \right] dt.
$$

Воспользуемся тригонометрическим тождеством для косинуса разности:
$$
\cos(\alpha - \beta) = \cos \alpha \cos \beta + \sin \alpha \sin \beta.
$$
Тогда выражение принимает вид:
\begin{equation}
    f(x) = \frac{1}{2l} \int_{-l}^{l} f(t) dt + \frac{1}{l} \sum_{k=1}^{\infty} \int_{-l}^{l} f(t) \cos\left( \frac{k\pi (t-x)}{l} \right) dt.
\end{equation}

\subsubsection{Предельный переход при $l \to \infty$}
Рассмотрим поведение каждого слагаемого при расширении интервала.

\textbf{1. Первое слагаемое:}
Так как по условию $f \in L_1(\mathbb{R})$, то интеграл $\int_{-\infty}^{+\infty} |f(t)| dt$ сходится к конечному числу. Следовательно, существует константа $C$, такая что:
$$
\left| \int_{-l}^{l} f(t) dt \right| \le \int_{-\infty}^{+\infty} |f(t)| dt = C.
$$
При делении на $2l$ (где $l \to \infty$), получаем:
$$
\frac{1}{2l} \int_{-l}^{l} f(t) dt \xrightarrow{l \to \infty} 0.
$$

\textbf{2. Второе слагаемое (сумма):}
Введем обозначения для дискретизации переменной:
$$
\lambda_k = \frac{k\pi}{l}, \quad \Delta \lambda = \lambda_k - \lambda_{k-1} = \frac{\pi}{l}.
$$
Выразим множитель перед интегралом: $\frac{1}{l} = \frac{\Delta \lambda}{\pi}$.
Тогда сумма приобретает вид интегральной суммы Римана:
$$
\frac{1}{\pi} \sum_{k=1}^{\infty} \Delta \lambda \underbrace{\left( \int_{-l}^{l} f(t) \cos(\lambda_k(t-x)) dt \right)}_{\text{Значение некоторой функции от } \lambda_k}.
$$

При $l \to \infty$ шаг разбиения $\Delta \lambda \to 0$, а пределы внутреннего интегрирования становятся бесконечными ($[-l, l] \to [-\infty, +\infty]$). Сумма Римана переходит в интеграл по переменной $\lambda$ от $0$ до $+\infty$.

Таким образом, получаем \textbf{интегральную формулу Фурье}:
\begin{equation}
    f(x) = \frac{1}{\pi} \int_{0}^{+\infty} d\lambda \int_{-\infty}^{+\infty} f(t) \cos(\lambda(t-x)) dt.
\end{equation}

Это выражение можно переписать в виде, аналогичном ряду Фурье:
\begin{equation}
    f(x) = \int_{0}^{+\infty} \left( a(\lambda) \cos(\lambda x) + b(\lambda) \sin(\lambda x) \right) d\lambda,
\end{equation}
где коэффициенты (теперь зависящие от непрерывного параметра $\lambda$) равны:
$$
a(\lambda) = \frac{1}{\pi} \int_{-\infty}^{+\infty} f(t) \cos(\lambda t) dt, \quad
b(\lambda) = \frac{1}{\pi} \int_{-\infty}^{+\infty} f(t) \sin(\lambda t) dt.
$$

\subsection{ Интеграл Фурье в комплексной форме}

Пусть $f \in L_1(\mathbb{R}) \cap C(\mathbb{R})$ и выполнено условие Дини.
Возьмем полученную формулу:
$$
f(x) = \frac{1}{\pi} \int_{0}^{+\infty} d\lambda \int_{-\infty}^{+\infty} f(t) \cos(\lambda(t-x)) dt.
$$

Заметим, что подынтегральная функция $\Psi(\lambda) = \int_{-\infty}^{+\infty} f(t) \cos(\lambda(t-x)) dt$ является \textbf{четной} относительно $\lambda$ (так как $\cos(-\alpha) = \cos(\alpha)$).
Интеграл от четной функции по $[0, +\infty)$ равен половине интеграла по всей оси $(-\infty, +\infty)$:
$$
f(x) = \frac{1}{2\pi} \int_{-\infty}^{+\infty} d\lambda \int_{-\infty}^{+\infty} f(t) \cos(\lambda(t-x)) dt.
$$

Рассмотрим теперь "мнимую" добавку. Интеграл
$$
\int_{-\infty}^{+\infty} f(t) \sin(\lambda(t-x)) dt
$$
является \textbf{нечетной} функцией от $\lambda$ (так как $\sin(-\alpha) = -\sin(\alpha)$). Интеграл от нечетной функции в симметричных пределах (в смысле главного значения) равен нулю:
$$
0 = \frac{1}{2\pi} \int_{-\infty}^{+\infty} d\lambda \int_{-\infty}^{+\infty} f(t) \sin(\lambda(t-x)) dt.
$$
Умножим этот ноль на $-i$ и добавим к выражению для $f(x)$:
$$
f(x) = \frac{1}{2\pi} \int_{-\infty}^{+\infty} d\lambda \int_{-\infty}^{+\infty} f(t) \left[ \cos(\lambda(t-x)) - i \sin(\lambda(t-x)) \right] dt.
$$
Используя формулу Эйлера $e^{-i\varphi} = \cos \varphi - i \sin \varphi$, получаем:
$$
f(x) = \frac{1}{2\pi} \int_{-\infty}^{+\infty} d\lambda \int_{-\infty}^{+\infty} f(t) e^{-i\lambda(t-x)} dt.
$$
Раскроем экспоненту $e^{-i\lambda(t-x)} = e^{-i\lambda t} \cdot e^{i\lambda x}$ и вынесем $e^{i\lambda x}$ из внутреннего интеграла (так как он не зависит от $t$):
\begin{equation}
    f(x) = \frac{1}{2\pi} \int_{-\infty}^{+\infty} d\lambda \, e^{i\lambda x} \underbrace{\int_{-\infty}^{+\infty} f(t) e^{-i\lambda t} dt}_{\text{Преобразование Фурье}}.
\end{equation}

\subsubsection{Определение преобразования Фурье}
Существуют различные соглашения о нормировочных коэффициентах.

\textbf{1. Несимметричный вариант:}
$$
F(\lambda) = \int_{-\infty}^{+\infty} f(x) e^{-i\lambda x} dx \quad \text{(прямое преобразование)}
$$
$$
f(x) = \frac{1}{2\pi} \int_{-\infty}^{+\infty} F(\lambda) e^{i\lambda x} d\lambda \quad \text{(обратное преобразование)}
$$

\textbf{2. Симметричный вариант:}
$$
F(\lambda) = \frac{1}{\sqrt{2\pi}} \int_{-\infty}^{+\infty} f(x) e^{-i\lambda x} dx
$$
$$
f(x) = \frac{1}{\sqrt{2\pi}} \int_{-\infty}^{+\infty} F(\lambda) e^{i\lambda x} d\lambda
$$

\textbf{Теорема (без доказательства):}
Пусть $f \in L_1(\mathbb{R})$. Если $\int_{-\infty}^{+\infty} f(t) e^{-i\lambda t} dt = 0$ для всех $\lambda$, то $f(x) = 0$ для почти всех $x \in \mathbb{R}$.

\subsection{4. Примеры}

\subsubsection{1. Экспоненциальное убывание}
Дано: $f(x) = e^{-\gamma|x|}$, где $\gamma > 0$.
Найдем $F(\lambda)$:
$$
F(\lambda) = \int_{-\infty}^{+\infty} e^{-\gamma|x|} e^{-i\lambda x} dx = \int_{-\infty}^{0} e^{\gamma x} e^{-i\lambda x} dx + \int_{0}^{+\infty} e^{-\gamma x} e^{-i\lambda x} dx.
$$
Вычислим интегралы:
$$
= \int_{-\infty}^{0} e^{(\gamma - i\lambda)x} dx + \int_{0}^{+\infty} e^{-(\gamma + i\lambda)x} dx =
\left. \frac{e^{(\gamma - i\lambda)x}}{\gamma - i\lambda} \right|_{-\infty}^0 - \left. \frac{e^{-(\gamma + i\lambda)x}}{\gamma + i\lambda} \right|_0^{+\infty}.
$$
На пределах бесконечности экспоненты зануляются (так как $\gamma > 0$), в нуле они равны 1:
$$
= \frac{1}{\gamma - i\lambda} + \frac{1}{\gamma + i\lambda} = \frac{\gamma + i\lambda + \gamma - i\lambda}{\gamma^2 + \lambda^2} = \frac{2\gamma}{\lambda^2 + \gamma^2}.
$$

\begin{center}
\begin{tikzpicture}
    \begin{axis}[
        width=8cm, height=5cm,
        axis lines=middle,
        xlabel=$x$, ylabel=$f(x)$,
        ymin=0, ymax=1.2,
        xmin=-3, xmax=3,
        ticks=none
    ]
    \addplot[blue, thick, domain=-3:3, samples=100] {exp(-abs(x))};
    \node[above] at (axis cs:0.5, 0.8) {$e^{-\gamma|x|}$};
    \end{axis}
\end{tikzpicture}
\quad
\begin{tikzpicture}
    \begin{axis}[
        width=8cm, height=5cm,
        axis lines=middle,
        xlabel=$\lambda$, ylabel=$F(\lambda)$,
        ymin=0, ymax=2.2,
        xmin=-3, xmax=3,
        ticks=none
    ]
    \addplot[red, thick, domain=-3:3, samples=100] {2/(x^2+1)};
    \node[above] at (axis cs:1, 1) {$\frac{2\gamma}{\lambda^2+\gamma^2}$};
    \end{axis}
\end{tikzpicture}
\end{center}

\subsubsection{2. Прямоугольный импульс}
Дано: $f(x) = \theta(a - |x|) = \begin{cases} 1, & |x| \le a \\ 0, & |x| > a \end{cases}$.
$$
F(\lambda) = \int_{-a}^{a} 1 \cdot e^{-i\lambda x} dx = \left. \frac{e^{-i\lambda x}}{-i\lambda} \right|_{-a}^{a} = \frac{e^{-i\lambda a} - e^{i\lambda a}}{-i\lambda} = \frac{e^{i\lambda a} - e^{-i\lambda a}}{i\lambda}.
$$
Так как $\frac{e^{ix} - e^{-ix}}{2i} = \sin x$, то:
$$
F(\lambda) = \frac{2\sin(\lambda a)}{\lambda}.
$$

\begin{center}
\begin{tikzpicture}
    \begin{axis}[
        width=8cm, height=4cm,
        axis lines=middle,
        xlabel=$x$, ylabel=$f(x)$,
        ymin=0, ymax=1.5,
        xmin=-3, xmax=3,
        xtick={-1.5, 1.5},
        xticklabels={$-a$, $a$}
    ]
    \addplot[blue, thick] coordinates {(-3,0) (-1.5,0) (-1.5,1) (1.5,1) (1.5,0) (3,0)};
    \end{axis}
\end{tikzpicture}
\quad
\begin{tikzpicture}
    \begin{axis}[
        width=8cm, height=4cm,
        axis lines=middle,
        xlabel=$\lambda$, ylabel=$F(\lambda)$,
        ymin=-0.5, ymax=2.5,
        xmin=-10, xmax=10,
        ticks=none
    ]
    \addplot[red, thick, domain=-10:10, samples=200] {2*sin(deg(x))/x};
    \node[right] at (axis cs:2, 1.5) {$\frac{2\sin(\lambda a)}{\lambda}$};
    \end{axis}
\end{tikzpicture}
\end{center}

\subsubsection*{3. Распределение Коши $\leftrightarrow$ Экспонента}
Функция $f(x) = \frac{1}{x^2 + a^2}$ переходит в $F(\lambda) = \frac{\pi}{a} e^{-a|\lambda|}$.
(Для графиков выбрано $a=1$).

\begin{center}
% График f(x)
\begin{tikzpicture}
    \begin{axis}[
        width=7cm, height=5cm,
        axis lines=middle,
        xlabel=$x$, ylabel=$f(x)$,
        ymin=0, ymax=1.2,
        xmin=-4, xmax=4,
        ticks=none
    ]
    % f(x) = 1/(x^2+1)
    \addplot[blue, thick, domain=-4:4, samples=100] {1/(x^2+1)};
    \node[above right] at (axis cs:0.5, 0.8) {$\frac{1}{x^2+a^2}$};
    \end{axis}
\end{tikzpicture}
\quad
% График F(lambda)
\begin{tikzpicture}
    \begin{axis}[
        width=7cm, height=5cm,
        axis lines=middle,
        xlabel=$\lambda$, ylabel=$F(\lambda)$,
        ymin=0, ymax=3.5, % pi approx 3.14
        xmin=-4, xmax=4,
        ticks=none
    ]
    % F(lambda) = pi * exp(-|x|)
    \addplot[red, thick, domain=-4:4, samples=100] {3.14159*exp(-abs(x))};
    \node[right] at (axis cs:0.5, 2) {$\frac{\pi}{a}e^{-a|\lambda|}$};
    \end{axis}
\end{tikzpicture}
\end{center}

% --- Пример 4: Гауссиана ---
\subsubsection*{4. Гауссов колокол $\leftrightarrow$ Гауссов колокол}
Функция $f(x) = e^{-ax^2}$ переходит в $F(\lambda) = \sqrt{\frac{\pi}{a}} e^{-\frac{\lambda^2}{4a}}$.
Преобразование Фурье сохраняет форму гауссианы (но меняет ширину и высоту).

\begin{center}
% График f(x)
\begin{tikzpicture}
    \begin{axis}[
        width=7cm, height=5cm,
        axis lines=middle,
        xlabel=$x$, ylabel=$f(x)$,
        ymin=0, ymax=1.2,
        xmin=-3, xmax=3,
        ticks=none
    ]
    % f(x) = exp(-x^2) (a=1)
    \addplot[blue, thick, domain=-3:3, samples=100] {exp(-x^2)};
    \node[right] at (axis cs:0.5, 0.7) {$e^{-ax^2}$};
    \end{axis}
\end{tikzpicture}
\quad
% График F(lambda)
\begin{tikzpicture}
    \begin{axis}[
        width=7cm, height=5cm,
        axis lines=middle,
        xlabel=$\lambda$, ylabel=$F(\lambda)$,
        ymin=0, ymax=2.0, % sqrt(pi) approx 1.77
        xmin=-6, xmax=6,
        ticks=none
    ]
    % F(lambda) = sqrt(pi)*exp(-x^2/4) (a=1)
    % Образ шире оригинала при a=1
    \addplot[red, thick, domain=-6:6, samples=100] {1.77*exp(-x^2/4)};
    \node[right] at (axis cs:1, 1.2) {$\sqrt{\frac{\pi}{a}} e^{-\frac{\lambda^2}{4a}}$};
    \end{axis}
\end{tikzpicture}
\end{center}

\chapter{Билет 13}
\section{ Аналог признака Дини для интеграла Фурье}

\subsection{Теорема}
Пусть выполняются следующие условия:

\begin{enumerate}
    \item Функция $f$ абсолютно интегрируема и непрерывна на всей числовой оси:
    $$ f \in L_1(\mathbb{R}) \cap C(\mathbb{R}) $$
    \item В фиксированной точке $x$ выполнено условие Дини:
    $$ \exists \delta > 0 : \quad \int_{-\delta}^{\delta} \left| \frac{f(x+t) - f(x)}{t} \right| \, dt < \infty $$
\end{enumerate}

\noindent \textbf{Тогда} в точке $x$ справедливо представление функции интегралом Фурье:
$$ f(x) = \frac{1}{\pi} \int_{0}^{+\infty} d\lambda \int_{-\infty}^{+\infty} f(t) \cos(\lambda(x-t)) \, dt $$

\subsection{Доказательство}

\subsubsection*{1. Рассмотрение урезанного интеграла и смена порядка интегрирования}
Рассмотрим интеграл до конечного числа $A$, обозначим его $J(A)$:
$$ J(A) = \frac{1}{\pi} \int_{0}^{A} d\lambda \int_{-\infty}^{+\infty} f(t) \cos(\lambda(x-t)) \, dt $$

Рассмотрим внутренний интеграл $\int_{-\infty}^{+\infty} f(t) \cos(\lambda(x-t)) \, dt$. Он сходится \textbf{равномерно} по параметру $\lambda \in [0, A]$.

\noindent \textit{Обоснование:}
Так как $\left| f(t) \cos(\dots) \right| \le |f(t)| \cdot 1 = |f(t)|$, и известно, что интеграл $\int_{-\infty}^{+\infty} |f(t)| \, dt$ сходится (по условию $f \in L_1(\mathbb{R})$), то по признаку Вейерштрасса сходимость исходного интеграла равномерная.

Из равномерной сходимости следует возможность поменять порядок интегрирования:
$$ J(A) = \frac{1}{\pi} \int_{-\infty}^{+\infty} f(t) \, dt \int_{0}^{A} \cos(\lambda(x-t)) \, d\lambda $$

\subsubsection*{2. Вычисление внутреннего интеграла}
Вычислим интеграл по $d\lambda$:
$$ \int_{0}^{A} \cos(\lambda(x-t)) \, d\lambda = \left. \frac{\sin(\lambda(x-t))}{x-t} \right|_{\lambda=0}^{\lambda=A} = \frac{\sin(A(x-t))}{x-t} $$
Заметим, что $\frac{\sin(A(x-t))}{x-t} = \frac{\sin(A(t-x))}{t-x}$ (так как синус и знаменатель меняют знак одновременно).

Подставим полученное выражение обратно в $J(A)$:
$$ J(A) = \frac{1}{\pi} \int_{-\infty}^{+\infty} f(t) \frac{\sin(A(t-x))}{t-x} \, dt $$

Сделаем замену переменной:
$$ \left[ \begin{gathered} u = t-x \\ t = x+u \end{gathered} \right] $$
В конспекте переменная интегрирования после замены вновь обозначена через $t$ (или $u$, но вернемся к $t$ для единообразия с интегралом Дирихле). Получаем:
$$ J(A) = \frac{1}{\pi} \int_{-\infty}^{+\infty} f(x+t) \frac{\sin(At)}{t} \, dt $$

\subsubsection*{3. Оценка разности}
Мы хотим доказать, что $J(A) \to f(x)$ при $A \to +\infty$.
Рассмотрим разность $J(A) - f(x)$.

Воспользуемся значением интеграла Дирихле (сходится условно):
$$ \int_{-\infty}^{+\infty} \frac{\sin(At)}{t} \, dt = \pi $$
Следовательно, умножив на константу $f(x)$, можем записать:
$$ f(x) = f(x) \cdot \frac{1}{\pi} \int_{-\infty}^{+\infty} \frac{\sin(At)}{t} \, dt = \frac{1}{\pi} \int_{-\infty}^{+\infty} f(x) \frac{\sin(At)}{t} \, dt $$

Тогда разность принимает вид:
$$ J(A) - f(x) = \frac{1}{\pi} \int_{-\infty}^{+\infty} \left( f(x+t) - f(x) \right) \frac{\sin(At)}{t} \, dt $$

\subsubsection*{4. Разбиение интеграла}
Для любого $\varepsilon > 0$ существует достаточно большое число $N_\varepsilon$ (граница "хвостов"). Разобьем интеграл на две части: внешнюю (хвосты) и внутреннюю (центр).

\begin{center}
\begin{tikzpicture}
    % Axis
    \draw[->] (-6,0) -- (6,0) node[right] {$t$};
    
    % Points
    \draw[thick] (-3,0.1) -- (-3,-0.1) node[below] {$-N_\varepsilon$};
    \draw[thick] (3,0.1) -- (3,-0.1) node[below] {$N_\varepsilon$};
    \draw[thick] (0,0.1) -- (0,-0.1) node[below] {$0$};
    
    % Braces
    \draw[decorate,decoration={brace,amplitude=10pt,mirror},thick] (-3,-0.5) -- (3,-0.5) node[midway,yshift=-20pt] {Центральная часть $I_2$};
    \draw[decorate,decoration={brace,amplitude=5pt},thick] (-6,0.5) -- (-3,0.5) node[midway,yshift=15pt] {Хвост $I_1$};
    \draw[decorate,decoration={brace,amplitude=5pt},thick] (3,0.5) -- (6,0.5) node[midway,yshift=15pt] {Хвост $I_1$};
\end{tikzpicture}
\end{center}

$$ |J(A) - f(x)| \le \underbrace{\left| \frac{1}{\pi} \int_{|t| \ge N_\varepsilon} (f(x+t) - f(x)) \frac{\sin(At)}{t} \, dt \right|}_{I_1} + \underbrace{\left| \frac{1}{\pi} \int_{-N_\varepsilon}^{N_\varepsilon} \frac{f(x+t) - f(x)}{t} \sin(At) \, dt \right|}_{I_2} $$

\paragraph{Оценка $I_1$ (Хвосты):}
Так как $f \in L_1(\mathbb{R})$, то интеграл от модуля функции по удаленным областям стремится к нулю.
При достаточно больших $|t|$ (т.е. $|t| \ge N_\varepsilon$) и с учетом того, что $\left| \frac{\sin(At)}{t} \right| \le \frac{1}{N_\varepsilon}$, вклад от хвостов можно сделать сколь угодно малым.
Более строго: так как $f \in L_1$, то $\int_{|t| \ge N_\varepsilon} |f(x+t)| \, dt \to 0$ при $N_\varepsilon \to \infty$.
Следовательно, можно выбрать $N_\varepsilon$ так, чтобы:
$$ I_1 < \frac{\varepsilon}{2} $$

\paragraph{Оценка $I_2$ (Центральная часть):}
Введем функцию $g(t)$ на отрезке $[-N_\varepsilon, N_\varepsilon]$:
$$ g(t) = \begin{cases} \frac{f(x+t) - f(x)}{t}, & t \in [-N_\varepsilon, N_\varepsilon] \\ 0, & \text{иначе} \end{cases} $$

Проверим, является ли $g(t)$ интегрируемой на $[-N_\varepsilon, N_\varepsilon]$ (то есть $g \in L_1([-N_\varepsilon, N_\varepsilon])$):
\begin{enumerate}
    \item В окрестностях нуля ($t \to 0$): интегрируемость обеспечена \textbf{условием Дини} ($\checkmark$).
    \item Вне окрестности нуля (на $[\delta, N_\varepsilon]$): значение $t$ отделено от $0$, а $f(x+t)$ непрерывна, значит $g(t)$ непрерывна и интегрируема.
\end{enumerate}
Следовательно, $g(t)$ — суммируемая функция. Тогда $I_2$ можно переписать как:
$$ I_2 = \left| \frac{1}{\pi} \int_{-N_\varepsilon}^{N_\varepsilon} g(t) \sin(At) \, dt \right| $$
По \textbf{лемме Римана} (об осцилляции) для $L_1$-функции:
$$ \int_{-N_\varepsilon}^{N_\varepsilon} g(t) \sin(At) \, dt \xrightarrow{A \to +\infty} 0 $$

Значит, существует $B > 0$, такое что для всех $A > B$:
$$ I_2 < \frac{\varepsilon}{2} $$

\subsubsection*{5. Итог}
Для любого $\varepsilon > 0$ мы выбрали $N_\varepsilon$ (для оценки хвостов), а затем нашли $B$ (для оценки центральной части), так что для всех $A > B$:
$$ |J(A) - f(x)| \le I_1 + I_2 < \frac{\varepsilon}{2} + \frac{\varepsilon}{2} = \varepsilon $$
Это и означает, что:
$$ J(A) \xrightarrow{A \to +\infty} f(x) $$
Теорема доказана. $\square$


\chapter{Билет 14}
\section*{Билет 14. Свойства преобразования Фурье}

Преобразование Фурье рассматривается как линейный оператор, действующий из пространства суммируемых функций $L_1(\mathbb{R})$ в пространство ограниченных функций $L_\infty(\mathbb{R})$.

\textbf{Определение.} Преобразованием Фурье функции $f(x)$ называется функция $F[f](\lambda)$, определяемая интегралом:
\[
F[f](\lambda) = \int_{-\infty}^{+\infty} f(x) e^{-i\lambda x} \, dx.
\]

\hrulefill

\subsection{1. Линейность и непрерывность}

\textbf{Теорема.} Оператор преобразования Фурье линеен и непрерывен. Если последовательность функций $f_n$ сходится к $f$ в пространстве $L_1(\mathbb{R})$ (то есть $\|f_n - f\|_{L_1} \to 0$ при $n \to \infty$), то последовательность образов $F[f_n]$ сходится к $F[f]$ равномерно на всей оси $\mathbb{R}$.

\textbf{Доказательство:}

Рассмотрим модуль разности образов:
\[
\left| F[f_n](\lambda) - F[f](\lambda) \right| = \left| \int_{-\infty}^{+\infty} (f_n(x) - f(x)) e^{-i\lambda x} \, dx \right|.
\]
Внесем модуль под знак интеграла. Так как $|e^{-i\lambda x}| = 1$ для любого вещественного $\lambda$, имеем:
\[
\left| \int_{-\infty}^{+\infty} (f_n(x) - f(x)) e^{-i\lambda x} \, dx \right| \le \int_{-\infty}^{+\infty} |f_n(x) - f(x)| \cdot 1 \, dx = \| f_n - f \|_{L_1(\mathbb{R})}.
\]
Таким образом, мы получили оценку, справедливую для всех $\lambda$:
\[
\sup_{\lambda \in \mathbb{R}} \left| F[f_n](\lambda) - F[f](\lambda) \right| \le \| f_n - f \|_{L_1(\mathbb{R})}.
\]
Правая часть стремится к нулю при $n \to \infty$ и \textbf{не зависит от $\lambda$}. Это означает, что сходимость равномерная.

\textbf{Следствие.} Справедлива оценка нормы оператора:
\[
\| F[f] \|_{C(\mathbb{R})} \le \| f \|_{L_1(\mathbb{R})}.
\]
Это доказывает, что преобразование Фурье — ограниченный оператор. \hfill $\square$

\subsection{2. Лемма Римана-Лебега}

\textbf{Лемма.} Преобразование Фурье переводит пространство $L_1(\mathbb{R})$ в подпространство $C_0(\mathbb{R})$ непрерывных функций, стремящихся к нулю на бесконечности:
\[
F[L_1(\mathbb{R})] \subset C_0(\mathbb{R}) = \left\{ \varphi \in C(\mathbb{R}) \mid \varphi(\lambda) \to 0 \text{ при } \lambda \to \infty \right\}.
\]

\textbf{Доказательство:}

Доказательство проводится в три этапа.

1. \textbf{Проверка для характеристической функции интервала.}
Пусть $f(x) = \chi_{[a, b]}(x)$ (равна 1 на отрезке и 0 вне его).
\[
F[\chi_{[a, b]}](\lambda) = \int_{a}^{b} e^{-i\lambda x} \, dx = \left. \frac{e^{-i\lambda x}}{-i\lambda} \right|_{a}^{b} = \frac{e^{-i\lambda b} - e^{-i\lambda a}}{-i\lambda}.
\]
При $\lambda \to \infty$ числитель ограничен (по модулю $\le 2$), а знаменатель растет. Следовательно, предел равен 0.

2. \textbf{Плотность ступенчатых функций.}
Линейные комбинации характеристических функций интервалов (ступенчатые функции) плотны в пространстве $L_1(\mathbb{R})$. Пусть $f \in L_1(\mathbb{R})$. Тогда существует последовательность ступенчатых функций $\{f_n\}$, сходящаяся к $f$ по норме $L_1$.

\begin{center}
\begin{tikzpicture}[scale=1]
    \draw[->] (-0.5,0) -- (6,0) node[right] {$x$};
    \draw[->] (0,-0.5) -- (0,3) node[above] {$y$};
    
    % Гладкая функция f(x)
    \draw[thick, blue] plot [smooth] coordinates {(0.5,0.5) (1.5,2.5) (3,1) (4.5,1.5) (5.5,0.2)};
    \node[blue] at (5.7, 0.5) {$f(x)$};
    
    % Ступенчатая аппроксимация f_n(x)
    \draw[red, thick] (0.5,0.5) -- (1,0.5) -- (1,1.5) -- (1.5,1.5) -- (1.5,2.5) -- (2,2.5) -- (2,1.75) -- (2.5,1.75) -- (2.5,1) -- (3,1) -- (3,1.25) -- (3.5,1.25) -- (3.5,1.5) -- (4,1.5) -- (4,1.2) -- (4.5,1.2) -- (4.5,0.8) -- (5,0.8) -- (5,0.2) -- (5.5,0.2);
    \node[red] at (3.5, 2.2) {$f_n(x)$};
    
    \node at (3, -0.7) {\textit{Аппроксимация функции ступенчатыми}};
\end{tikzpicture}
\end{center}

3. \textbf{Предельный переход.}
Из свойства 1 следует, что $F[f_n] \rightrightarrows F[f]$ (равномерно).
Поскольку $F[f_n]$ — конечная линейная комбинация образов характеристических функций, то каждая $F[f_n]$ непрерывна и стремится к нулю на бесконечности.
Равномерный предел непрерывных функций есть непрерывная функция, поэтому $F[f] \in C(\mathbb{R})$.

Рассмотрим предел на бесконечности. Так как сходимость равномерная, можно поменять порядок пределов:
\[
\lim_{\lambda \to \infty} F[f](\lambda) = \lim_{\lambda \to \infty} \lim_{n \to \infty} F[f_n](\lambda) = \lim_{n \to \infty} \underbrace{\lim_{\lambda \to \infty} F[f_n](\lambda)}_{=0} = 0.
\]
\hfill $\square$

\subsection{3. Связь производной и умножения (дифференцирование оригинала)}

\textbf{Свойство.} Пусть $f \in C^1(\mathbb{R}) \cap L_1(\mathbb{R})$ (непрерывно дифференцируема) и её производная $f' \in L_1(\mathbb{R})$. Тогда:
\[
F[f'](\lambda) = i\lambda F[f](\lambda).
\]

\textbf{Доказательство:}

Для начала заметим, что если $f \in L_1$ и $f' \in L_1$, то существуют пределы $\lim_{x \to \pm \infty} f(x)$. Так как $f \in L_1$, эти пределы обязаны быть равны нулю (иначе интеграл от функции расходился бы). Следовательно:
\[
f(x) \to 0 \quad \text{при } x \to \pm \infty.
\]

Рассмотрим интеграл для преобразования производной и применим интегрирование по частям:
\[
F[f'](\lambda) = \int_{-\infty}^{+\infty} f'(x) e^{-i\lambda x} \, dx.
\]
Введем обозначения для интегрирования по частям:
\[
u = e^{-i\lambda x} \implies du = -i\lambda e^{-i\lambda x} \, dx, \qquad dv = f'(x) \, dx \implies v = f(x).
\]
Подставляем:
\[
\int_{-\infty}^{+\infty} f'(x) e^{-i\lambda x} \, dx = \underbrace{\left. f(x) e^{-i\lambda x} \right|_{-\infty}^{+\infty}}_{=0} - \int_{-\infty}^{+\infty} f(x) (-i\lambda) e^{-i\lambda x} \, dx.
\]
Внеинтегральный член зануляется из-за убывания $f(x)$. Остается:
\[
= i\lambda \int_{-\infty}^{+\infty} f(x) e^{-i\lambda x} \, dx = i\lambda F[f](\lambda).
\]
\hfill $\square$

\subsection{4. Обобщение на высшие порядки}

\textbf{Свойство.} Если $f \in C^l(\mathbb{R})$ и все производные до порядка $l$ включительно суммируемы ($f^{(k)} \in L_1(\mathbb{R}), k=0, \dots, l$), то:
\[
F[f^{(l)}](\lambda) = (i\lambda)^l F[f](\lambda).
\]

\textbf{Следствие (скорость убывания спектра).} Выразим $F[f](\lambda)$ из формулы выше:
\[
F[f](\lambda) = \frac{F[f^{(l)}](\lambda)}{(i\lambda)^l}.
\]
По лемме Римана-Лебега числитель $F[f^{(l)}](\lambda) \to 0$ при $\lambda \to \infty$. Следовательно:
\[
F[f](\lambda) = o\left( \frac{1}{\lambda^l} \right) \quad \text{при } \lambda \to \infty.
\]
Это означает, что чем более гладкой является функция, тем быстрее убывает её преобразование Фурье.

\subsection{5. Дифференцирование образа}

\textbf{Свойство.} Если $f(x) \in L_1(\mathbb{R})$ и $x \cdot f(x) \in L_1(\mathbb{R})$, то преобразование Фурье дифференцируемо, и выполняется равенство:
\[
(F[f])'(\lambda) = F[-ix f(x)].
\]

\textbf{Доказательство:}

Запишем выражение для $F[f](\lambda)$ и продифференцируем его по параметру $\lambda$. Дифференцирование под знаком интеграла допустимо, так как интеграл сходится равномерно (в силу условия $x f(x) \in L_1$).
\[
F(\lambda) = \int_{-\infty}^{+\infty} f(x) e^{-i\lambda x} \, dx,
\]
\[
F'(\lambda) = \frac{d}{d\lambda} \int_{-\infty}^{+\infty} f(x) e^{-i\lambda x} \, dx = \int_{-\infty}^{+\infty} f(x) \frac{\partial}{\partial \lambda} \left( e^{-i\lambda x} \right) \, dx.
\]
Так как $\frac{\partial}{\partial \lambda} (e^{-i\lambda x}) = -ix e^{-i\lambda x}$, получаем:
\[
F'(\lambda) = \int_{-\infty}^{+\infty} \left[ -ix f(x) \right] e^{-i\lambda x} \, dx = F[-ix f(x)](\lambda).
\]
\hfill $\square$

\subsection{6. Обобщение дифференцирования образа}

\textbf{Свойство.} Если $f(x), x f(x), \dots, x^k f(x) \in L_1(\mathbb{R})$, то образ $F[f]$ принадлежит классу $C^k(\mathbb{R})$, и для производной $k$-го порядка справедлива формула:
\[
F^{(k)}[f] = F\left[ (-ix)^k f(x) \right].
\]

\subsection{7. Класс Шварца}

\textbf{Определение.} Пространством (классом) Шварца $S(\mathbb{R})$ называется пространство бесконечно дифференцируемых функций, которые убывают на бесконечности быстрее любой степени полинома вместе со всеми своими производными.
Формально:
\[
S = \left\{ f \in C^\infty(\mathbb{R}) \mid \forall p, q \ge 0 \quad \exists C_{p,q} > 0 : \sup_{x \in \mathbb{R}} \left| x^p f^{(q)}(x) \right| < C_{p,q} \right\}.
\]

\textbf{Важное свойство.} Преобразование Фурье является биекцией пространства Шварца на себя:
\[
F: S \to S \quad \text{— биекция}.
\]
Это свойство делает класс Шварца крайне удобным для работы с преобразованием Фурье, так как операции дифференцирования и умножения на полином (которые соответствуют друг другу при преобразовании) не выводят функцию за пределы этого класса.
\chapter{Билет 15}
\section{Формулировка теоремы Фубини для интеграла Лебега. Существование свёртки суммируемых функций. Теорема
о свёртке.}
\textbf{Тема:} Формулировка теоремы Фубини для интеграла Лебега. Существование свёртки суммируемых функций. Теорема о свёртке.

\bigskip

\textbf{Опр. (Свёртка)}

Пусть $f_1, f_2$ — функции на $\mathbb{R}$. Свёрткой функций $f_1$ и $f_2$ (обозначается $(f_1 * f_2)(x)$) называют функцию, определяемую интегралом:
$$
(f_1 * f_2)(x) := \int_{-\infty}^{+\infty} f_1(t) f_2(x-t) \, dt
$$

\hrulefill

\bigskip

\subsection{Теорема (Фубини)}

Пусть $(X, \mu)$ и $(Y, \nu)$ — пространства с мерой ($\mu$ — мера на $X$, $\nu$ — мера на $Y$). На декартовом произведении $X \times Y$ задана мера $\lambda = \mu \times \nu$.
Пусть функция $f(x, y)$ суммируема на произведении, то есть $f \in L_1(X \times Y, \lambda)$.

\textbf{Утверждение:}
\begin{enumerate}
    \item Для почти всех $x \in X$ (по мере $\mu$) существует интеграл по $y$: $\displaystyle \int_Y f(x, y) \, d\nu(y)$.
    \item Для почти всех $y \in Y$ (по мере $\nu$) существует интеграл по $x$: $\displaystyle \int_X f(x, y) \, d\mu(x)$.
    \item Повторные интегралы равны двойному:
    $$
    \int_{X \times Y} f(x, y) \, d\lambda(x, y) = \int_X \left( \int_Y f(x, y) \, d\nu(y) \right) d\mu(x) = \int_Y \left( \int_X f(x, y) \, d\mu(x) \right) d\nu(y)
    $$
\end{enumerate}

\bigskip

\subsection{Утверждение (Теорема Тонелли)}

Если функция $f(x, y)$ измерима и существует (конечен) хотя бы один из повторных интегралов от модуля $|f|$:
$$
\int_X \left[ \int_Y |f(x, y)| \, d\nu \right] d\mu < \infty \quad \text{или} \quad \int_Y \left( \int_X |f(x, y)| \, d\mu(x) \right) d\nu(y) < \infty,
$$
то $f \in L_1(X \times Y, \lambda)$ и к ней применима теорема Фубини.

\hrulefill

\bigskip

\subsection{Теорема (О существовании свёртки)}

Если $f_1, f_2 \in L_1(\mathbb{R})$ (абсолютно интегрируемы на прямой), то:
\begin{enumerate}
    \item Для почти всех $x \in \mathbb{R}$ интеграл свёртки существует.
    \item $f_1 * f_2 \in L_1(\mathbb{R})$.
\end{enumerate}

\textbf{Доказательство:}

Так как $f_1, f_2 \in L_1$, то $\int_{\mathbb{R}} |f_1(t)|dt < \infty$ и $\int_{\mathbb{R}} |f_2(s)|ds < \infty$.
Рассмотрим произведение норм:
$$
\exists \, I = \int_{\mathbb{R}} dt |f_1(t)| \cdot \int_{\mathbb{R}} |f_2(s)| \, ds < \infty.
$$
Сделаем замену переменной во внутреннем интеграле: пусть $s = x - t$. Тогда при фиксированном $t$ имеем $ds = dx$ (или $dx = ds$). Пределы интегрирования остаются $\mathbb{R}$.
$$
I = \int_{\mathbb{R}} dt |f_1(t)| \int_{\mathbb{R}} |f_2(x-t)| \, dx.
$$
Это повторный интеграл от неотрицательной функции $|f_1(t)f_2(x-t)|$. Так как он конечен, по \textbf{теореме Тонелли} функция суммируема на плоскости, и мы можем поменять порядок интегрирования:
$$
\int_{\mathbb{R}} dx \int_{\mathbb{R}} |f_1(t) f_2(x-t)| \, dt = I < \infty.
$$
Раз внешний интеграл по $x$ от выражения $\left( \int |f_1(t) f_2(x-t)| dt \right)$ сходится, то подынтегральная функция конечна для почти всех $x \in \mathbb{R}$.
$$
\Rightarrow \text{при п.в. } x \in \mathbb{R} \quad \exists \int_{\mathbb{R}} f_1(t) f_2(x-t) \, dt.
$$
Это доказывает п. 1.

Для п. 2 оценим $L_1$-норму свёртки:
$$
\|f_1 * f_2\|_{L_1} = \int_{\mathbb{R}} \left| \int_{\mathbb{R}} f_1(t) f_2(x-t) \, dt \right| dx \le \int_{\mathbb{R}} dx \int_{\mathbb{R}} |f_1(t) f_2(x-t)| \, dt.
$$
Как показано выше (через $I$), этот интеграл равен $\|f_1\|_{L_1} \cdot \|f_2\|_{L_1} < \infty$.
$$
\Rightarrow f_1 * f_2 \in L_1(\mathbb{R}). \quad \blacksquare
$$




\subsection{Теорема (О свёртке и преобразовании Фурье)}

Если $f_1, f_2 \in L_1(\mathbb{R})$, то $F[f_1 * f_2] = F[f_1] \cdot F[f_2]$.

\textbf{Доказательство:}

Запишем определение преобразования Фурье от свёртки:
$$
F[f_1 * f_2](\lambda) = \int_{-\infty}^{+\infty} (f_1 * f_2)(x) e^{-i\lambda x} \, dx = \int_{-\infty}^{+\infty} dx \left( \int_{-\infty}^{+\infty} f_1(t) f_2(x-t) \, dt \right) e^{-i\lambda x}.
$$
Так как подынтегральная функция (как доказано выше) суммируема, то по \textbf{т. Фубини} можно поменять порядок интегрирования:
$$
= \int_{-\infty}^{+\infty} dt \, f_1(t) \int_{-\infty}^{+\infty} f_2(x-t) e^{-i\lambda x} \, dx.
$$
Во внутреннем интеграле сделаем замену переменных:
$$
\left[ \begin{gathered} y = x - t \\ x = y + t \\ dx = dy \end{gathered} \right]
$$
Тогда экспонента распадается: $e^{-i\lambda x} = e^{-i\lambda(y+t)} = e^{-i\lambda y} \cdot e^{-i\lambda t}$.
$$
= \int_{-\infty}^{+\infty} dt \, f_1(t) \int_{-\infty}^{+\infty} f_2(y) e^{-i\lambda y} e^{-i\lambda t} \, dy.
$$
Выносим множитель $e^{-i\lambda t}$, зависящий только от $t$, из внутреннего интеграла, а интеграл по $y$ (зависящий только от $\lambda$) — из внешнего:
$$
= \left( \int_{-\infty}^{+\infty} f_2(y) e^{-i\lambda y} \, dy \right) \cdot \left( \int_{-\infty}^{+\infty} f_1(t) e^{-i\lambda t} \, dt \right).
$$
Первый множитель — это $F[f_2](\lambda)$, второй — $F[f_1](\lambda)$.
$$
= F[f_2](\lambda) \cdot F[f_1](\lambda) = F[f_1](\lambda) \cdot F[f_2](\lambda). \quad \blacksquare
$$
\end{document}

